\section{Mathematical results so far}
\label{sec:mathematical-summary}

\noindent\textbf{6 September: explicit spectral and umbral maps, and the original family's parameter identifications.}
For the literal common lattice $\mathcal I=\{\sqrt2(a,a,a):a\in D_4\}$,
write $a=a_0+v$, $t^2=|v|^2$, $Q=a_0^2+t^2$ and
$P=a_0(a_0^2-3t^2)$. Its complete cubic image is
\[
 \tau(\mathcal I)=2\sqrt2\bigl(18\mathbb Z\cup(6\mathbb Z+2)
                                      \cup(6\mathbb Z-2)\bigr).
\]
Three explicit integral affine lines give a section of this value map.
The scalar-retaining Albert matrix $X=\mu I+Y(a)$, with cyclic
off-diagonal entries $\sqrt2a$, has invariants and labelled spectrum
\[
\begin{aligned}
 (\operatorname{tr}X,\operatorname{tr}X^2,\det X)
   &=(3\mu,3\mu^2+12Q,\mu^3-6\mu Q+4\sqrt2P),\\
 \operatorname{Spec}Y
   &=(2\sqrt2a_0,-\sqrt2a_0-\sqrt6t,-\sqrt2a_0+\sqrt6t).
\end{aligned}
\]
For $t>0$, the explicit labelled eigenprojector satisfies
$(\Pi_1)_{12}=-1/6-\sqrt3(v/t)/6$; thus the frame and spectrum
recover $a$, not only its norm and cubic. The full discriminant is
$864t^2(3a_0^2-t^2)^2$. The full Cartan family retains the scalar
$\mu$, radius $r$, and angle $\theta$, with $r^2=12Q$ and
$r^3\cos3\theta=24\sqrt3P$ on this subfamily.

For the same marked root system $A_5^4D_4$, the
Cheng--Duncan--Harvey prescription gives
$S^X=(5I_{12}+J)S_6$, where $J_{rs}=\mathbf1_{s=5r\ (12)}$ and
$S_{6,r}=\sum_{k\equiv r\ (12)}kq^{k^2/24}$.
The coefficient map to the retained Fricke function is
\[
 T_6=\frac{\eta(\tau)^5\eta(3\tau)}{\eta(2\tau)\eta(6\tau)^5},
 \qquad f^X=-q\frac{d}{dq}\log T_6,
 \qquad T_6\!\left(-\frac1{6\tau}\right)=\frac{72}{T_6(\tau)}.
\]
All twelve theta indices, the multiplier determinant, the original
glue-group action and its twisted shadows are computed in the proofs.
This is an explicit use and comparison of established umbral machinery,
not a claim to a new moonshine module.

In the original analytic family $\beta_c=\beta^{\rm part}+c$,
$\operatorname{Im}c<-M$, the actual global identifications are exactly
$X_c\simeq X_{c'}\Longleftrightarrow c'-c\in2\mathbb Z$.
The middle lattices agree for every integer shift; the multiplicity-four
filling retains an order-two obstruction to odd shifts. Even shifts have
explicit finite-filling and toric cusp maps and a proved global additive
gluing correction. The resulting parameter quotient is
$w=e^{-\pi ic}$, $0<|w|<e^{-\pi M}$.
This classifies the retained analytic family without using the claimed
identification of its total space with $S^6$.


\paragraph{6 September: cubic arithmetic, integral lattice maps, and the canonical ring.}
For the retained $A_5^4\perp D_4$ Niemeier lattice $N$ and its
original index-six Leech neighbour $\Lambda$, fix the Cayley--Dickson
frame $(1,i,j,k,\ell,i\ell,j\ell,k\ell)$ used in the complete cubic
proof. Put $\tau(q)=\operatorname{Re}((q_1q_2)q_3)$ and
$F=\tau\circ\mathcal L$ on the original Niemeier span,
and $\mathcal V_F(L)=\operatorname{span}_{\mathbb Z}\{F(y):y\in L\}$.
The exact additive groups are
\[
\begin{aligned}
 \mathcal V_F(N)&=\tfrac1{27}\mathbb Z\oplus
       \tfrac{\sqrt2}{54}\mathbb Z\oplus\tfrac{\sqrt3}{18}\mathbb Z,\\
 \mathcal V_F(\Lambda)&=\tfrac1{27}\mathbb Z\oplus
       \tfrac{\sqrt2}{216}\mathbb Z\oplus\tfrac{\sqrt3}{54}\mathbb Z,
 \qquad \mathcal V_F(\Lambda)/\mathcal V_F(N)\cong\mathbb Z/12.
\end{aligned}
\]
The Albert determinant is $\lambda^3-\lambda(y,y)+2F(y)$.
The quotient concerns the displayed value-generated groups in this fixed
frame, not an inclusion $N\subset\Lambda$. The actual residue
$r=\rho(361)$ cycles three Leech neighbours $\Lambda_j=r^j\Lambda$.
Their pairwise intersection is one integral lattice $J$, with
$\det J=81$ and $[\Lambda_j:J]=9$. The fixed fourth isotropic plane
in $J^\vee/J$ gives $M$ with exactly $72$ roots of type $A_2^{12}$,
$M\cap\Lambda_j=J$, and $[M:J]=9$.
\[
 \mathcal V_F(J)=\mathcal V_F(M)
 =\tfrac1{27}\mathbb Z\oplus\tfrac{\sqrt2}{216}\mathbb Z
                         \oplus\tfrac{\sqrt3}{18}\mathbb Z.
\]
The successive value-group inclusions have matrices
$\operatorname{diag}(1,4,1)$ and $\operatorname{diag}(1,1,3)$,
giving quotients $\mathbb Z/4$ and $\mathbb Z/3$.
Both containments and their reverse generator certificates are calculated
from the full cubic polarization, not inferred from lattice indices.

Wilson's construction \cite{esn:Wilson2009}, transported by the explicit
isometry $\iota(x)=x/\sqrt2$, gives $\Lambda_{\rm cyc}$ in the
original triality metric. Its cyclic fixed Gram matrix is $3G_{E_8}$;
its orthogonal complement is explicitly isometric to $A_2\otimes E_8$,
and their sum has index $3^8$, with quotient $\mathbb F_3^8$.
The directions remain $c(q_1,q_2,q_3)=(q_3,q_1,q_2)$ and
$c_{\rm W}=c^{-1}$. Its literal intersections are
\[
 \Lambda_{\rm cyc}\cap\mathcal L(N)
 =\Lambda_{\rm cyc}\cap\mathcal L(\Lambda)
 =\Delta:=\{\sqrt2(a,a,a):a\in D_{\mathbb H}\},
\]
where $D_{\mathbb H}$ is the even-coordinate-sum integer lattice in
the first four stated Cayley coordinates. Thus $\Delta$ has rank $4$,
Gram matrix $6G_{D_4}$, determinant $5184$, minimum squared norm $12$,
and value-generated subgroup $4\sqrt2\mathbb Z$ for $\tau|_\Delta$, or
$8\sqrt2\mathbb Z$ for $2\tau|_\Delta$. This is not the actual image:
$12\sqrt2$ is not attained by $\tau$ on $\Delta$.
The preceding proofs retain the triality products and spectral boundary
graph, the flag's degree-$+1$ collapse, and the retained cusp class's
nontrivial three-plane reduction; these exact maps do not supply an
interacting quantum mass-gap result.

For the original analytic fibration $f:X\to\mathbb P^1$, the full
canonical-divisor proof gives
$\omega_X\simeq\mathcal O_X(-2S_2)$ and the graded anticanonical ring
$\bigoplus_{m\ge0}H^0(X,\omega_X^{-m})\cong\mathbb C[U,V]$,
with $\deg U=1$, $\deg V=2$. The complete system $|-2K_X|$
recovers $f$ with the specified projective coordinates. This calculation
uses the original analytic charts and finite fillings, not an
identification of $X$ with $S^6$, and draws no sphere-topology conclusion.

The retained constant $\beta_c=\beta^{\rm part}+c$ now gives a proper
holomorphic deformation over $\operatorname{Im}c<-M$, locally on
relatively compact parameter discs. On a fixed regular fibre set
$T=\operatorname{Im}\tau$, $m=\operatorname{Im}\mu$,
$D=\operatorname{Im}\beta_{c_*}-6m^2/T<0$ and $\delta=c-c_*$.
In the plus-graph convention its exact Beltrami matrix is
\[
 \nu_\delta=\frac{\delta}{2iD+\delta}
                 \begin{pmatrix}0&0\\-m/T&1\end{pmatrix}.
\]
The connecting Kodaira--Spencer class is
$-[\partial_{\zeta_2}\otimes
(d\bar\zeta_2-(m/T)d\bar\zeta_1)/(2iD)]\ne0$.
The intrinsic pencil recovers $f$ also over the dual numbers; its three
fixed critical values prohibit an infinitesimal projective base change.
Consequently this deformation of the whole $X_c$ is nontrivial,
not only its marked regular-fibre deformation. The proof uses no sphere
recognition or CDP conclusion.



\paragraph{5 September: exact corrections to two S6 proof steps.}
The Cartier duality isomorphism in Theorem~\ref{thm:cartier-duality}
is induced by the analytic trace with $K_W^\bullet=\omega_W[2]$;
outer identities of five-term sequences alone do not define it.
For the holomorphic overlap translation by $\alpha(s)\Pi(s)v_j$,
the real marking gives
$a_j=\operatorname{Re}\alpha\,v_j+\operatorname{Im}\alpha\,J_sv_j$,
not multiplication of a real vector by the complex scalar $\alpha$.
The proved covariance $J_{g_js}A_j=A_jJ_s$ and $A_jv_j=v_j$
give $a_j(g_js)=A_ja_j(s)-v_j/m_j$. On the full clockwise
covering circle, $\Delta a_j=-v_j$ and the mixed class of
$h_j^*\Omega_E-\Omega_E$ is $-\iota_{v_j}E$.
Thus Lemma~\ref{lem:picard-finite-norm} still gives
$N_{T_j}d_j=\iota_{v_j}E$, and
Proposition~\ref{prop:geometric-pi1} retains its based meridian relation.
The complete proofs now include the base two-form term and its zero
restriction to the circle; no coefficient or orientation is suppressed.

\paragraph{5 September: exact period constants and cusp edge maps.}
For the cusp-regular particular solution define
\[
 D_{\rm part}=\operatorname{Im}\beta^{\rm part}
       -\frac{6(\operatorname{Im}\mu)^2}{\operatorname{Im}\tau},
 \qquad M=\max_{B\setminus\{p_0\}}D_{\rm part}.
\]
Proposition~\ref{prop:period-exact-admissible-constants} proves that
$M$ is finite and attained, and that $\beta=\beta^{\rm part}+c_0$
gives an everywhere invertible marked real period map exactly for
$\operatorname{Im}c_0<-M$. The cusp-free affine solution
$\widetilde\beta=\beta+i\,t\circ\pi$ has
$\widetilde D=D+\operatorname{Re}t\to+\infty$ along the positive
real base ray. The actual cusp extension excludes it.
For $C_v=E_0\cap E_v$, the cusp identification is
$\iota_v=\Psi_{B_0v}|_{C_v}:C_v\to C_{-v}$.
Using the same character $w=\chi^{(a,0)}$ with $av=0$ on both curves,
\[
 \iota_v^*w_{-v}=e^{2\pi i\,aC(0)B_0v}w_v.
\]
For $(v,a)=(e_1,(0,1)),(e_2,(1,0)),(e_1-e_2,(1,1))$, the multipliers
are respectively $e^{-2\pi i\mu_0}$, $e^{12\pi i\mu_0}$ and
$e^{-2\pi i(7\mu_0+b_0)}$, with $\mu_0=\mu(0)$ and $b_0=b(0)$.
Equation~\eqref{eq:cusp-exact-opposite-edge} retains the full $C(0)$;
Lemma~\ref{lem:toric-chart-control} uses restricted charts and the
closed-polydisc condition $\theta_i^\sigma(y)\ge0$ for torus points.
These refinements and the bounded finite-filling audit do not supply
an additional sphere-identification or CDP theorem.

\paragraph{5 September: the complete two-cross scalar--Gauss domain.}
Section~\ref{q444:sgd:section} retains the literal matrices, constants,
and all three full scalar and Gauss residuals. Put $\delta=\sqrt{2/3}$ and
\[
 \mathcal D=\{(u,v_0,y)\in\mathbb R^3:
 u>0,\ d=uv_0-\delta^2-y^2>0,\ c=4v_0-5u>0\}.
\]
Theorem~\ref{sgd:parametrization} proves a bijection from
$\mathcal D\times\{1,-1\}^2\times U(1)^2$ to all classified-frame
solutions, with $t=72d/c$, $e=\epsilon\sqrt{2+u+v_0}$,
$k=\sigma\sqrt{d/(-\det T)}$ and $\omega=e/k$, where
\[
 \begin{gathered}
 M=\begin{pmatrix}u&\delta+iy\\\delta-iy&v_0\end{pmatrix},\quad
 Q=\operatorname{diag}(1,2),\\
 T=\frac t4\operatorname{diag}(2/3,-5/6)
   -\frac12(QM+MQ).
 \end{gathered}
\]
The rank-one matrices $A_\pm=M\pm kT$ have $h_\pm=(A_\pm)_{11}>0$;
the exact vectors and phase inverse are
\[
 \begin{gathered}
 a_\pm=\begin{pmatrix}\sqrt{h_\pm}\\(A_\pm)_{21}/\sqrt{h_\pm}\end{pmatrix},\qquad
 p_\pm=\frac{(v\mp iw)_1}{\sqrt{h_\pm}},\\
 v=\frac{p_+a_++p_-a_-}{2},\quad
 w=\frac{i(p_+a_+-p_-a_-)}2.
 \end{gathered}
\]
Theorem~\ref{sgd:literal-addon:quotient} proves that the full literal
domain is the quotient of these parameters and $\zeta\in U(1)$ by
the sole free involution
$(p_+,p_-,\zeta)\mapsto(-p_+,-p_-,-\zeta)$.
The exact inverse frame recovers every literal temporal and spatial entry.
The smooth phase coordinates are
$(\eta,\alpha,\beta)=(\zeta^2,p_+/\zeta,p_-/\zeta)$.
Theorems~\ref{sgd:smooth-addon:literal-diffeomorphism}
and~\ref{sgd:smooth-addon:coupling} give
\[
 \mathcal S_{\rm sg}\cong\{1,-1\}^2\times(0,\infty)
                         \times\mathbb R^2\times\mathbb T^3,
 \qquad t=\tau=g^2,
\]
with coupling coordinates and inverse
\[
 v_0=\frac{\sqrt5}{2}\sqrt{\delta^2+y^2}+e^s,\qquad
 u=\frac{4\tau v_0+72(\delta^2+y^2)}{72v_0+5\tau},
 \quad (\tau,y,s)\in(0,\infty)\times\mathbb R^2.
\]
The literal manifold has four embedded six-dimensional components;
each positive-coupling level has four embedded five-dimensional
components $\mathbb R^2\times\mathbb T^3$
(Corollary~\ref{sgd:smooth-addon:levels}).
For each original $\tau=g^2>0$ the choice
\[
 \begin{gathered}
 \rho=\frac{\tau+\sqrt{\tau^2+4608\delta^2}}{96\delta}
       >\frac1{\sqrt2},\\
 M_\rho=\delta\begin{pmatrix}\rho&1\\1&2\rho\end{pmatrix},\qquad
 t=\frac{24\delta(2\rho^2-1)}\rho=\tau
 \end{gathered}
\]
proves nonemptiness with both temporal and frequency signs
(Corollary~\ref{sgd:every-coupling}). Both scalar equations and every
Gauss entry vanish, while Theorem~\ref{q444:twocross:thm:exclusion}
proves that the two spatial residuals cannot vanish simultaneously
on this entire class.
The quantitative spatial bound in
Proposition~\ref{q444:twocross:quantitative} therefore has an explicitly
nonempty domain. Existence on the unrestricted remaining exceptional
set stays \Open; no quantum gap claim follows.

\paragraph{5 September higher-rung and selector continuation.}
Section~\ref{sec:h24-continuation} gives the complete higher-wrap portion
of the expanded root draft.
The explicit octonionic slice has determinant
$\lambda^3-\lambda\sum_i|q_i|^2+2\operatorname{Re}((q_1q_2)q_3)$
and positive-cone compactification $(B^{25},S^{24})$.
The flag $F_4/\operatorname{Spin}(8)$ has Poincar\'e polynomial
$1+2t^8+2t^{16}+t^{24}$ and Euler class $6\beta_1^2\beta_2$.
For the retained peripheral character,
$(\chi(C_0),\chi(C_1),\chi(C'_2))=(-1,4,-3)$ in $\mathbb Z/12$.
The exact suspension is $E^6(a\alpha)=2a\nu_9$ in
$\pi_{12}(S^9)\cong\mathbb Z/24$; the particular cusp coefficient is
now $a=0$. Theorem~\ref{h24:thm:cusp-g-zero} constructs the complete
framed filling inside one original chart times an interval, retaining
the full period functions. Thus $g\simeq *$ and $\Theta_g\chi=0$.
The retained-angle map is different: Theorem~\ref{h24:ret:main}
proves $[p_3q_2]=\eta_4\eta_5\ne0$ in $\pi_6(S^4)$.
Its full normally four-framed inverse torus has spin form
$q(x,y)=x+xy+y$ over $\mathbb F_2$ and Arf invariant one,
so the original conical quotient $q_2$ is nonnull. The exact
homomorphism $[x]\mapsto[\eta_2\wedge x]$ sends this value to
$12\nu_6\ne0$ in $\pi_9(S^6)\cong\mathbb Z/24$; its
threefold suspension is $12\nu_9$ in the same group receiving
$E^6[g]=0$. Composing with the original $\chi$ gives peripheral
values $(12,0,12)$ and kernel $\chi^{-1}(2\mathbb Z/12)$.
The full class and exact order of $q_2$ are not computed by this
order-two quotient alone.
The arithmetic addition in
Section~\ref{sec:fifteen-selector-jordan} maps the complete fifteen
$p=1201$ selector rays into the real symmetric corner of
$\operatorname{Herm}_3(\mathbb H)$ and proves both linear span and
Jordan closure have dimension three. These additions retain the global
higher-completion question explicitly.

\paragraph{The modulo-$840$ lattice, full Weyl quotient, and retained counts.}
Section~\ref{esn:sec:shared} proves
\[
 R_{\mathrm{ES}}=A_5^4\perp D_4\subset N,\qquad
 [N:R_{\mathrm{ES}}]=72,\qquad \det N=1,\qquad
 \#\operatorname{Roots}(N)=144,
\]
\[
 O(N)=W(R_{\mathrm{ES}})\rtimes L(G),\qquad
 G\cong GL_2(\mathbb F_3)\cong Q_8\rtimes S_3.
\]
The full quotient and residue-position theorems
\ref{esn:fa:thm:full-quotient} and~\ref{esn:fa:thm:residue-weyl}
give $\rho(S_{840})\subset W(R_{\mathrm{ES}})$ and
$L_g\rho(r)=\rho(r)L_g$ in the original six-coordinate fibre order.
The abstract quotient is credited to \cite[Table~3]{esn:CDH2014};
the explicit code isomorphism, lattice lift and marked $D_4$ weight
comparison are proved here. The shell maps retain the actual prime:
$A_3(2521)=0$, $A_3(3361)=6$, with unordered counts $0,3$.
For a finite set $P$ of primes reducing into $S_{840}=\langle289\rangle$,
the count pushforward gives
$Q_{R,P}=6W_{R,P}-M_{R,P}u\in A(S_{840})\subset N$;
Theorem~\ref{esn:thm:sample-fourier} gives the exact Fourier inverse
for the residue-count function $w_{R,P}$.

This is the short route through the current mathematics.  It states what is
proved in this workbench and exhibits the load-bearing formulas.  Both the
intrinsic nearby-cycle/Leray certificate and the independent
collapse/Mayer--Vietoris certificate are now complete for the candidate.

\begin{enumerate}[leftmargin=1.8em,itemsep=.65em]
\item \Verified\ \textbf{The analytic threefold exists in the reconstructed
coordinate model.}
The descended modular periods give
\[
 \Pi=[Z\mid I_2],\qquad
 Z=\begin{pmatrix}6\mu&\tau\\\beta&\mu\end{pmatrix},
\]
with
\[
 \det_{\mathbb R}\Pi
 =-\det\operatorname{Im}Z
 =\operatorname{Im}\tau
 \left(\operatorname{Im}\beta-
 \frac{6(\operatorname{Im}\mu)^2}{\operatorname{Im}\tau}\right)<0.
\]
The Cayley coordinates
\(s_j(z)=(z-z_j)/(z-\bar z_j)\), the two elliptic orientations, and the
modular deck normalization are fixed simultaneously:
\[
 g_1'(z_1)=e^{-2\pi i/3},\qquad g_2'(z_2)=e^{-2\pi i/4}=-i,
 \qquad
 s_1\circ g_1=e^{-2\pi i/3}s_1,\quad s_2\circ g_2=-is_2.
\]
The period-deck homomorphism $\varphi$ is defined by
$\tau\circ g=\varphi(g)\circ\tau$.  Its cusp image is proved in a fixed
covering component, not by an untyped single-valuedness assertion: choose
\[
 U_0\subset\pi^{-1}(D_0^*),\qquad g_0U_0=U_0,\qquad
 \tau(U_0)\subset K\subset j^{-1}(\{|w|>R\}).
\]
For $P=\varphi(g_0)$ one has $P(K)=K$.  If $c$ is the boundary cusp of
$K$, choose $\gamma_0(\infty)=c$ and put
\[
 \tau'=\gamma_0^{-1}\tau|_{U_0},\qquad
 P'=\gamma_0^{-1}P\gamma_0,\qquad
 \tau'\circ g_0=P'\circ\tau'.
\]
If $P'=1$, then $\tau'$ descends to
$\bar\tau':D_0^*\to\mathfrak h$ and
$q'=\exp(2\pi i\bar\tau')$ simultaneously has a positive-order zero at the
cusp and a single-valued holomorphic logarithm, an impossible winding
combination.  Thus $P'$ and hence $P$ are nontrivial parabolic.

The auxiliary $\gamma_0$-frame is then discarded and the proof returns to
$\tau(z_1)=\rho$, $\varphi(g_1)=S_1$.  For the trace-$+2$ lift, with
primitive fixed vector $(p,r)^{\mathsf t}$, $\gcd(p,r)=1$, and $k\ne0$,
\[
 \widetilde P=I+k
 \begin{pmatrix}-pr&p^2\\-r^2&pr\end{pmatrix},
\]
and
\[
 \widetilde S_1=\begin{pmatrix}1&-1\\1&0\end{pmatrix},\qquad
 \operatorname{tr}(\widetilde S_1^{-1}\widetilde P^{-1})
 =1+k(p^2-pr+r^2)=0.
\]
Thus projective order two of $S_1^{-1}P^{-1}$ gives
$-k(p^2-pr+r^2)=1$, so $k=-1$ and
$p^2-pr+r^2=1$; conjugating the whole deck representation by a power of
$S_1$ leaves $S_1$ fixed and gives the final normalization below.  Rename
the resulting period map $\tau$ and define anew, in this cusp paragraph,
\[
 q=e^{2\pi i\tau}=t_cu_1(t_c),\qquad u_1(0)\ne0.
\]
For
\(S_1(w)=(w-1)/w\), \(S(w)=-1/w\), and \(P(w)=w-1\),
\[
 \varphi(g_1)=S_1,\qquad \varphi(g_2)=S,\qquad
 \varphi(g_0)=P,\qquad S_1SP=1
 \quad\text{in }\operatorname{PSL}_2(\mathbb Z).
\]
At $z_2$, comparison of quadratic leading terms gives
\[
 \bigl(g_2'(z_2)\bigr)^2=S'(i)=-1;
\]
the simultaneous parabolic normalization selects $g_2'(z_2)=-i$.
The periodic toric cusp action is free and properly discontinuous; its
bounded-chart argument uses
\(\lvert B_t\bar\lambda\rvert\ge\frac12\lvert\bar\lambda\rvert\).
Exponentiation identifies the punctured toric quotient with the family of
compact two-dimensional complex tori, not with an abelian-surface family.
For the finite logarithmic actions the opposite translation signs are
retained:
\[
 v_1=\varepsilon,\qquad v_2=-\varepsilon',\qquad
 \gamma(v_1)=1,\qquad\gamma(v_2)=-1,
\]
and
\[
 (z,x)\longmapsto(g_jz,A_jx+v_j/m_j).
\]
These actions are free by the invariant covector criterion.  On an overlap,
the compatible logarithm is
\[
 \ell'(w)=\ell(g_j^{-1}w)-\frac{2\pi i}{m_j},
 \qquad
 h_j\widetilde g_j^{\log}h_j^{-1}=\widetilde g_j^0,
\]
so no branch sign is suppressed.  Their radial stalk maps are the printed
quotient maps,
and the three fillings glue to a compact connected smooth complex
threefold \(X\) with a proper surjective holomorphic map
\[
 f:X\longrightarrow\mathbb P^1.
\]
On the cusp cover the toric form
\[
 \Omega=\frac{dx_1\wedge dx_2\wedge dt}{x_1x_2}
       =\pm dz_0\wedge dz_1\wedge dz_2
\]
extends nowhere vanishingly and satisfies
\(\Psi_{\bar\lambda}^*\Omega=\Omega\).  This first proves
\(\omega_{N_0}\simeq\mathcal O_{N_0}\); only then does the open embedding
\(N_0\hookrightarrow X\) give the typed two-stage chain
\[
 \omega_X|_{N_0}\simeq\omega_{N_0}\simeq\mathcal O_{N_0}.
\]
The complete chain is
\cref{prop:period-input-audit,thm:candidate-cusp-geometry,prop:candidate-cusp-overlap,prop:candidate-finite-fillings,lem:finite-filling-stalk-map,thm:candidate-global-geometry}.

\item \Verified\ \textbf{The cusp fibre has the exact geometry needed in the
topology and conductor arguments.}
Its scheme-theoretic local equations are
\[
 z_0,\qquad z_0z_1,\qquad z_0z_1z_2.
\]
It is reduced, irreducible, nonnormal, and normal crossings (not simple
normal crossings).  Its normalization is \(\mathrm{dP}_6\), with opposite
sides of the anticanonical hexagon identified, and its double locus consists
of three rational curves through two triple points.

\item \Verified\ \textbf{The algebraic reduction is computed directly and
does not use the topology or CDP20.}
For a very general smooth fibre,
\[
 \operatorname{NS}(F_b)=\mathbb Z\eta,\qquad
 \eta=u\wedge w+6\gamma\wedge\delta,\qquad
 \eta\wedge\eta=12\gamma\wedge u\wedge w\wedge\delta.
\]
The inverse-dual matrix \(G_\eta\) associated with the Appell--Humbert form
has determinant
\[
 \det G_\eta=
 \frac23\operatorname{Im}\tau
 \left(\operatorname{Im}\beta-
 \frac{6(\operatorname{Im}\mu)^2}{\operatorname{Im}\tau}\right)<0,
\]
and the same inertia as that form, hence signature \((1,1)\) and
\(a(F_b)=0\).  The repaired
meromorphic-descent argument then gives
\[
 \mathcal M(X)=f^*\mathbb C(t),\qquad a(X)=1,
\]
and \(f\) is the algebraic reduction; see
\cref{prop:very-general-torus-audit,thm:algebraic-dimension-audit}.

\item \Verified\ \textbf{The intrinsic nearby-cycle/Leray route proves the
integral topology without using the cellular collapse.}
For the cusp filling, the integral nearby-cycle map satisfies
\[
 \operatorname{sp}^r:H^r(W;\mathbb Z)
   \xrightarrow{\ \sim\ }H^r(F;\mathbb Z)^{T_0}
 \qquad(0\le r\le4)
\]
as saturated lattices.  With
$q:=uw+6\gamma\delta\in\bigwedge^2V=H^2(F;\mathbb Z)$, the global sheaf calculation gives
\[
H^0(R^1)=\mathbb Z(12\gamma),\quad
H^0(R^2)=\mathbb Z(2q),\quad
H^0(R^3)=\mathbb Z(2\gamma uw).
\]
The canonical truncation triangle
\[
 R^0f_*\mathbb Z_X\longrightarrow\tau_{\le1}Rf_*\mathbb Z_X
 \longrightarrow R^1f_*\mathbb Z_X[-1]
 \xrightarrow{\kappa}R^0f_*\mathbb Z_X[1]
\]
identifies the explicitly computed Cech discrepancy with Leray \(d_2\).
For \((\ell_0,\ell_1,\ell_2)=(0,1,-1)\), hence \(p=-1\), orientation of
\(\omega\) may be chosen so that
\[
d_2(12\gamma)=\omega,\qquad
d_2(2q)=\omega[\delta],\qquad
d_2(2\gamma uw)=\omega[\gamma\delta].
\]
Therefore
\[
 H^1(X;\mathbb Z)=H^2(X;\mathbb Z)=H^3(X;\mathbb Z)=0,
 \qquad H_2(X;\mathbb Z)=0.
\]

\item \Verified\ \textbf{The based van Kampen maps and sphere recognition
close the topology theorem.}
Zero-section whiskers fix the common basepoint, the finite logarithmic
overlaps impose the primitive relations with one common sign, the cusp
inclusion is \(\Lambda\to\bar\Lambda\), and the cusp meridian bounds the
explicit toric zero-section disc with no translation drift.  Hence
\[
\pi_1(X)=
\langle c,x,y\mid c\text{ central},\ xy=1,\ x^3=c,\ y^4=c^{-1}\rangle
\simeq1.
\]
Poincar\'e duality, universal coefficients, Hurewicz, and the smooth
homotopy-sphere recognition theorem now give
\[
 X\cong_{\mathrm{diff}}S^6,\qquad
 \langle c_3(T_X),[X]\rangle=\chi_{\mathrm{top}}(X)=2.
\]
See
\cref{thm:p0-integral-specialization,prop:leray-e2,lem:first-transgression,lem:higher-transgressions,thm:h2-vanishing,prop:geometric-pi1,cor:candidate-sphere}.

\item \Refuted\ \textbf{The printed CDP dependency has two exact failures.}
For a reduced nonnormal fibre,
\[
0\longrightarrow N^*_{S/X}\longrightarrow
\ker(\Omega_X^1|_S\to\widetilde\Omega_S^1)
\longrightarrow\operatorname{tors}\Omega_S^1\longrightarrow0,
\]
so the vanishings used on p.~692 of \cite{CDP2020} do not control the
conductor-supported middle term.  A proper
nodal-cubic-times-\(\mathbb P^1\) family is a counterexample to that
inference.  Separately, a split extension satisfies
\[
 (K\rtimes Q)^{\mathrm{ab}}
 \simeq Q^{\mathrm{ab}}\oplus(K^{\mathrm{ab}})_Q,
\]
not the direct sum with all of \(K^{\mathrm{ab}}\) unless monodromy acts
trivially.

\item \Verified\ \textbf{The singular-fibre nonvanishing has a
topology-free $\operatorname{Pic}^0$ proof and three exact all-Pic
certificates.}
The invariant toric three-form first trivializes $\omega_{N_0}$, and the
open embedding then gives
$\omega_X|_{N_0}\simeq\omega_{N_0}\simeq\mathcal O_{N_0}$, as recorded
above.  Let
$\nu:\widetilde W\simeq\mathrm{dP}_6\to W$ be the normalization.  For
$L\in\operatorname{Pic}^0(X)$, naturality of $c_1$ and injectivity of
\[
 c_1:\operatorname{Pic}(\mathrm{dP}_6)\hookrightarrow
 H^2(\mathrm{dP}_6;\mathbb Z)
\]
give
$\nu^*((L^\vee\otimes\omega_X)|_W)\simeq\mathcal O_{\mathrm{dP}_6}$.
The conductor section, Ramis--Ruget duality on the compact Cartier fibre,
and top-degree analytic base change then prove
\[
 R^2f_*(T_X\otimes L)\ne0
 \qquad(L\in\operatorname{Pic}^0(X))
\]
without any global homology calculation.  For arbitrary $L$, the exact
trivial-pullback obstruction is the homomorphism
\[
 r:\operatorname{Pic}(X)\longrightarrow\operatorname{Pic}(\mathrm{dP}_6),
 \qquad r([L])=[\nu^*i^*L].
\]
For the candidate, the Leray edge image of $c_1(L)$ lies in
$E_\infty^{0,2}\subseteq\ker d_2^{0,2}$, but
$d_2^{0,2}(2q)=\omega[\delta]$ is an isomorphism.  Hence the edge section is
zero; proper base change and naturality give $c_1(L|_W)=0$, and injectivity
of $c_1$ on $\mathrm{dP}_6$ forces
\[
 \nu^*i^*L\simeq\mathcal O_{\mathrm{dP}_6},
 \qquad r([L])=0.
\]
This proves the conclusion for
every $L\in\operatorname{Pic}(X)$ without citing the packaged
$H^2(X;\mathbb Z)=0$ theorem, although it reuses the same integral Leray
transgression.

The second certificate uses no global Leray differential.  The exact
order-four filling lattice forces $c_1(L|_F)=2nq$.  Subtracting the invariant
vertical form on the pair-of-pants bundle gives a crossed homomorphism with
global-coinvariant boundary values.  The logarithmic translations and their
cyclic norm maps give
\[
 \bar d_1=4n[\delta],\qquad \bar d_2=-3n[\delta],
\]
while the flat zero-section cusp sweep bounds a literal toric solid torus,
so $\bar d_0=0$.  The based relation gives
$\bar d_1+\bar d_2+\bar d_0=0$, hence $n=0$.  Naturality and injectivity of
the local integral cusp specialization then give $c_1(L|_W)=0$ and
\[
 \nu^*i^*L\simeq\mathcal O_{\mathrm{dP}_6},
 \qquad r([L])=0.
\]
This is independent of global $d_2$, though not of local cusp
topology.  See
\cref{prop:pic0-pullback,prop:all-pic-restriction,%
prop:leray-edge-restriction,prop:direct-boundary-picard,%
cor:candidate-pic0-route,cor:candidate-conductor-actual}.

The precise limit of a cusp-only analytic replacement is also computed.
For $Q(v)=v_1^2-6v_2^2$, the locally finite toric divisor
$\mathcal D_Q=\sum_vQ(v)E_v$ has an explicit equivariant descent
$\mathcal L_Q\in\operatorname{Pic}(N_0)$ satisfying
\[
 c_1(\mathcal L_Q|_{F_t})=-2q,\qquad
 \bigl(\deg_{C_i}\nu^*(\mathcal L_Q|_W)\bigr)
 =(-12,2,0,-12,2,0).
\]
Thus $\operatorname{Pic}(N_0)\to\operatorname{Pic}(\mathrm{dP}_6)$ is
nonzero, so no cusp-local representability argument can prove the global
restriction theorem by itself.  The completed third certificate adds the two
finite fillings analytically and retains the prior coinvariant input
\[
 (\textstyle\bigwedge^2V)^G=\mathbb Zq,\qquad
 q=u\wedge w+6\gamma\wedge\delta,\qquad
 V_G\otimes\mathbb R=\mathbb R[\delta].
\]
Thus, for a global line bundle $L$,
$E=c_1(L|_F)=kq$.  The finite norm proof keeps every covering level
distinct.  On the vector cover
\[
 \widetilde J_j=\Delta_j\times\mathbb C^2
 \xrightarrow{\ r_j\ }J'_j
\]
define the vector-cover automorphisms
\[
 Q_{j,\log},Q_{j,\mathrm{lin}}\in\operatorname{Aut}(\widetilde J_j),
 \qquad
 Q_{j,\log}=\widetilde g_j^{\,v_j/m_j},\quad
 Q_{j,\mathrm{lin}}=\widetilde g_j^{\,0}.
\]
These are not the quotient automorphisms: on the vector cover,
\[
\begin{aligned}
 Q_{j,\log}\tau_\lambda Q_{j,\log}^{-1}
 &=\tau_{A_j\lambda},&
 (Q_{j,\log})^{m_j}&=\tau_{v_j},\\
 Q_{j,\mathrm{lin}}\tau_\lambda Q_{j,\mathrm{lin}}^{-1}
 &=\tau_{A_j\lambda},&
 (Q_{j,\mathrm{lin}})^{m_j}&=1.
\end{aligned}
\]
The sign in $(Q_{j,\log})^{m_j}=\tau_{v_j}$ is positive because
$m_j$ clockwise iterations of
$s_j\circ g_j=e^{-2\pi i/m_j}s_j$ produce translation by $+v_j$.
Let
\[
 q_{j,\log},q_{j,\mathrm{lin}}\in\operatorname{Aut}(J'_j)
\]
be their descents.  They have order dividing $m_j$ by these relations and
exactly $m_j$ because they cover $g_j$.  On
$J_j^{\prime *}=J'_j|_{\Delta_j^*}$ define
\[
\begin{aligned}
 p_{j,\log}&:J_j^{\prime *}\longrightarrow
 N_j^*=J_j^{\prime *}/\langle q_{j,\log}\rangle,\\
 p_{j,\mathrm{lin}}&:J_j^{\prime *}\longrightarrow
 J_j^*=J_j^{\prime *}/\langle q_{j,\mathrm{lin}}\rangle .
\end{aligned}
\]
The branch-independent translation and its descended biholomorphism are
typed by
\[
 h_j:J_j^{\prime *}\longrightarrow J_j^{\prime *},\qquad
 \bar h_j:N_j^*\longrightarrow J_j^*,
\]
and satisfy
\[
 h_j\circ q_{j,\log}\circ h_j^{-1}=q_{j,\mathrm{lin}},\qquad
 p_{j,\mathrm{lin}}\circ h_j=\bar h_j\circ p_{j,\log},
\]
and the filled line bundle gives the typed pullback comparison
\[
 p_{j,\log}^*(L|_{N_j^*})
 \simeq h_j^*p_{j,\mathrm{lin}}^*(L|_{J_j^*}).
\]

For the factor extension return to the filled vector cover:
\[
 P_{j,\log}:J'_j\longrightarrow N_j,\qquad
 \pi_j^{\rm vec}:=P_{j,\log}\circ r_j:
 \widetilde J_j\longrightarrow N_j,\qquad
 P_{j,\log}|_{J_j^{\prime *}}=p_{j,\log},
\]
\[
 \Gamma_j^{\log}
 =\langle\tau_\lambda,Q_{j,\log}:\lambda\in\Lambda\rangle.
\]
The pullback $(\pi_j^{\rm vec})^*(L|_{N_j})$ defines the filled factor
extension.  Over $\Delta_j^*$ the quotient square identifies its
commutator form with $E$; that integer-valued form is locally constant on
the connected disc, so it remains $E$ on the filled cover.  After
trivializing this pullback, choose normalized vertical lifts
\[
 U_\lambda^{(j)}U_\mu^{(j)}
 =\mathbf c\!\left(\frac12E(\lambda,\mu)\right)
 U_{\lambda+\mu}^{(j)},
\]
write
$\alpha\tau_\lambda\alpha^{-1}=\tau_{A_\alpha\lambda}$ and
$T_\alpha=(A_\alpha^{-1})^{\mathsf t}$, and define
$\widetilde d_j(\alpha)\in V_{\mathbb R}$ by
\[
 \widetilde G_\alpha U_{A_\alpha^{-1}\lambda}^{(j)}
 \widetilde G_\alpha^{-1}
 =\mathbf c\!\left(
 \langle\widetilde d_j(\alpha),\lambda\rangle\right)
 U_\lambda^{(j)}.
\]
It is a crossed homomorphism:
\[
 \widetilde d_j(\alpha\beta)
 =\widetilde d_j(\alpha)+T_\alpha\widetilde d_j(\beta).
\]
To compare the boundary factors, use the actual common cover and its
restricted lattice quotient
\[
 \widetilde J_j^*:=\Delta_j^*\times\mathbb C^2,
 \qquad
 r_j^*:=r_j|_{\widetilde J_j^*}:
 \widetilde J_j^*\longrightarrow J_j^{\prime *}.
\]
The two quotient maps from this common domain are
\[
\begin{aligned}
 \Phi_{j,\log}:\widetilde J_j^*&\longrightarrow N_j^*,&
 \Phi_{j,\log}&:=p_{j,\log}\circ r_j^*,\\
 \Phi_{j,\mathrm{glob}}:\widetilde J_j^*&\longrightarrow J_j^*,&
 \Phi_{j,\mathrm{glob}}&:=p_{j,\mathrm{lin}}\circ h_j\circ r_j^*
 =\bar h_j\circ\Phi_{j,\log}.
\end{aligned}
\]
Both maps are $\Gamma_j^{\log}$-invariant.  Translation invariance follows
from $r_j^*\circ\tau_\lambda=r_j^*$.  For every
$\widetilde x\in\widetilde J_j^*$, invariance under the nontrivial generator
is the point-valued calculation
\[
\begin{aligned}
 \Phi_{j,\log}\bigl(Q_{j,\log}(\widetilde x)\bigr)
 &=p_{j,\log}\bigl(q_{j,\log}(r_j^*(\widetilde x))\bigr)
 =\Phi_{j,\log}(\widetilde x),\\
 \Phi_{j,\mathrm{glob}}\bigl(Q_{j,\log}(\widetilde x)\bigr)
 &=p_{j,\mathrm{lin}}\bigl(
 h_j(q_{j,\log}(r_j^*(\widetilde x)))\bigr)\\
 &=p_{j,\mathrm{lin}}\bigl(
 q_{j,\mathrm{lin}}(h_j(r_j^*(\widetilde x)))\bigr)
 =\Phi_{j,\mathrm{glob}}(\widetilde x).
\end{aligned}
\]
Thus $\widetilde J_j^*$ is the same $\Gamma_j^{\log}$-cover of the two
biholomorphic quotients $N_j^*$ and $J_j^*$.  Pulling back the gluing
isomorphism gives
\[
 \Theta_j:
 \Phi_{j,\log}^*(L|_{N_j^*})
 \xrightarrow{\ \sim\ }
 \Phi_{j,\mathrm{glob}}^*(L|_{J_j^*}).
\]
It intertwines descent and hence identifies the pushed-out filled and global
$\Gamma_j^{\log}$-factor extensions while fixing every central
$\mathbf c(r)$, $r\in\mathbb R$.

The loop and its three successive lifts are fixed pointwise as follows.  Let
\[
 \gamma_j^J:[0,1]\longrightarrow J_j^*
\]
be the zero-section image of the unwhiskered clockwise boundary meridian,
so $\gamma_j^J(1)=\gamma_j^J(0)$.  Choose the global zero-section basepoint
$\widehat b\in J|_P$ and the fixed zero-section whisker
\[
 \omega_j:[0,1]\longrightarrow J|_P,\qquad
 \omega_j(0)=\widehat b,\qquad
 \omega_j(1)=\gamma_j^J(0),
\]
where the boundary copy of $J_j^*$ is identified with its image in $J|_P$.
The retained based total-space loop is the class
\[
 \widehat\rho_j
 =[\omega_j*\gamma_j^J*\omega_j^{-1}]
 \in\pi_1(J|_P,\widehat b).
\]
Choose $x_j\in J_j^{\prime *}$ such that
$p_{j,\mathrm{lin}}(h_j(x_j))=\gamma_j^J(0)$, put
$x_{j,\mathrm{lin}}=h_j(x_j)$, and let
$\widetilde\gamma_{j,\mathrm{lin}}:[0,1]\to J_j^{\prime *}$ be the
$p_{j,\mathrm{lin}}$-lift of $\gamma_j^J$ beginning at
$x_{j,\mathrm{lin}}$.  For the chosen clockwise generator,
\[
\begin{aligned}
 p_{j,\mathrm{lin}}\circ\widetilde\gamma_{j,\mathrm{lin}}
 &=\gamma_j^J,\\
 \widetilde\gamma_{j,\mathrm{lin}}(0)&=x_{j,\mathrm{lin}},\\
 \widetilde\gamma_{j,\mathrm{lin}}(1)
 &=q_{j,\mathrm{lin}}(x_{j,\mathrm{lin}}).
\end{aligned}
\]
Set
\[
 \gamma_j^N:=\bar h_j^{-1}\circ\gamma_j^J:[0,1]\longrightarrow N_j^*,
 \qquad
 \widetilde\gamma_{j,\log}:=h_j^{-1}\circ
 \widetilde\gamma_{j,\mathrm{lin}}:[0,1]\longrightarrow J_j^{\prime *}.
\]
Then the quotient square gives the point-valued identities
\[
\begin{aligned}
 p_{j,\log}\circ\widetilde\gamma_{j,\log}&=\gamma_j^N,\\
 \widetilde\gamma_{j,\log}(0)&=x_j,\\
 \widetilde\gamma_{j,\log}(1)&=q_{j,\log}(x_j).
\end{aligned}
\]
Choose $\widetilde x_j\in\widetilde J_j^*$ with
$r_j^*(\widetilde x_j)=x_j$, and let
$\widetilde\gamma_{j,\mathrm{vec}}:[0,1]\to\widetilde J_j^*$ be the
$r_j^*$-lift of $\widetilde\gamma_{j,\log}$ beginning at
$\widetilde x_j$.  On adjacent sheets choose compatible logarithm branches
\[
 \ell'(w)=\ell(g_j^{-1}w)-\frac{2\pi i}{m_j}.
\]
The point-valued vector-cover endpoint is
\[
\begin{aligned}
 r_j^*\circ\widetilde\gamma_{j,\mathrm{vec}}
 &=\widetilde\gamma_{j,\log},\\
 \widetilde\gamma_{j,\mathrm{vec}}(0)&=\widetilde x_j,\\
 \widetilde\gamma_{j,\mathrm{vec}}(1)
 &=Q_{j,\log}(\widetilde x_j).
\end{aligned}
\]
excluding both $Q_{j,\log}^{-1}(\widetilde x_j)$ and every
$\tau_\lambda Q_{j,\log}(\widetilde x_j)$ with $\lambda\ne0$.  Fibre loops
remain the pointwise deck paths from $\widetilde x_j$ to
$\tau_\lambda(\widetilde x_j)$, since $h_j$ is a fibre translation.

The factor of $L$ along the fixed lift of $\omega_j$ identifies the target
factor over the local boundary segment with the globally based factor.  Let
$G_{\rho_j}^{(j)}$ be the target lift of the local boundary deck element
whose whiskered image is the chosen global lift $G_{\rho_j}$.  It satisfies
\[
 G_{\rho_j}^{(j)}U_{A_j^{-1}\lambda}
 (G_{\rho_j}^{(j)})^{-1}
 =\mathbf c(\langle d_j,\lambda\rangle)U_\lambda,
 \qquad d_j=d(\rho_j).
\]
Consequently, for some $a_j\in V_{\mathbb R}$ and $b_j\in\mathbb R$,
\[
 \Theta_{j,*}(U_\lambda^{(j)})
 =\mathbf c(\langle a_j,\lambda\rangle)U_\lambda,\qquad
 \Theta_{j,*}(\widetilde G_{Q_{j,\log}})
 =\mathbf c(b_j)G_{\rho_j}^{(j)}.
\]
The central factor $\mathbf c(b_j)$ disappears under conjugation.  Therefore
\[
 \widetilde d_j(Q_{j,\log})
 =d_j+(T_j-I)a_j,\qquad
 \widetilde d_j(\tau_v)=\iota_vE.
\]
Here
$A_{Q_{j,\log}}=A_j$ and $T_{Q_{j,\log}}=T_j$.  Consequently the exact,
sign-sensitive norm chain is
\[
\begin{aligned}
 N_{T_j}d_j
 &=N_{T_j}\widetilde d_j(Q_{j,\log})
 =\widetilde d_j((Q_{j,\log})^{m_j})\\
 &=\widetilde d_j(\tau_{v_j})
 =\iota_{v_j}E.
\end{aligned}
\]
Finally, on a simply connected $U\subset\Delta_j^*$ the logarithm branch
$\ell$ gives the point-typed vector-cover lift
\[
 H_j^\ell:U\times\mathbb C^2\longrightarrow U\times\mathbb C^2,
 \qquad
 H_j^\ell(z,\zeta)
 =\left(z,\zeta+\frac{\ell(z)}{2\pi i}\Pi(z)v_j\right)
\]
of $h_j$.  A branch change $\ell\mapsto\ell+2\pi i n$ gives
\[
 H_j^{\ell+2\pi i n}=\tau_{nv_j}\circ H_j^\ell.
\]
At factor-extension level this gives
\[
 \Theta'_{j,*}
 =\operatorname{Ad}(U_{nv_j})\circ\Theta_{j,*}.
\]
Thus the comparison character changes by
$a_j\mapsto a_j+n\iota_{v_j}E$.  Since
\[
 T_j(\iota_{v_j}E)=\iota_{v_j}E,
\]
both displayed $\widetilde d_j$ and norm identities are unchanged; the
descended $h_j,\bar h_j$ and quotient square are branch independent.
The exact contractions are
\[
 \iota_{v_1}q=2\eta_1,\quad \iota_{v_2}q=-3\eta_2,\quad
 [\eta_1]=3[\delta],\quad [\eta_2]=2[\delta],
\]
and logarithms of the actual $\mathcal O^*$ factors give
\[
 \bar d_1=2k[\delta],\qquad
 \bar d_2=-\tfrac32k[\delta],\qquad
 \bar d_0=0,\qquad
 \bar d_1+\bar d_2+\bar d_0=0,
\]
where the cusp equality is an argument-principle calculation in a Stein
chart.  Hence $k=0$ and $E=0$.

The toric support-function calculation uses the explicitly typed unimodular
map
\[
 B_0:\bar\Lambda\simeq\mathbb Z^2
 \longrightarrow\Lambda_{\rm tor}\simeq\mathbb Z^2,
 \qquad
 [B_0]=\begin{pmatrix}0&1\\-1&0\end{pmatrix},
 \qquad \det B_0=1.
\]
For adjacent rays $\rho_v,\rho_w<\sigma=\{v,w,u\}$ the affine-open
directions are, exactly,
\[
 U_{\rho_v}=D(z_wz_u)\subset U_\sigma,\qquad
 U_{\rho_w}=D(z_vz_u)\subset U_\sigma,\qquad
 U_{\rho_v}\cap U_{\rho_w}=D(z_vz_wz_u)=U_0.
\]
If
$\mu^+_\lambda:\Psi_\lambda^*\mathcal O(D_a)\to\mathcal O(D_a)$
is multiplication by $\chi^{m_\lambda}$, the map needed on the negative
divisor is the inverse-dual, not the ordinary dual:
\[
 \mu^-_\lambda:=((\mu^+_\lambda)^{-1})^\vee:
 \Psi_\lambda^*\mathcal O(-D_a)\longrightarrow\mathcal O(-D_a),
 \qquad s\longmapsto\chi^{-m_\lambda}s.
\]
Now set, with its type retained,
\[
 M=L|_{N_0}\in\operatorname{Pic}(N_0),\qquad
 E_M=c_1(M|_{F_t})=E=0.
\]
For
$A=\nu^*(M|_W)$ and
$(\deg_{C_1}A,\ldots,\deg_{C_6}A)=(d_1,d_2,d_3,d_1,d_2,d_3)$,
the normalized degrees determine
\[
 [B]=\begin{pmatrix}d_2+d_3&d_3\\d_3&d_1+d_3\end{pmatrix},
 \qquad
 E_M((\lambda,0),(0,m))=-B(B_0\lambda,m).
\]
The antecedent $E_M=0$ and unimodularity of $B_0$ force $B=0$, hence all
six degrees vanish and
\[
 \nu^*(L|_W)=\nu^*(M|_W)\simeq\mathcal O_{E_0};
 \qquad
 r([L])=[\mathcal O_{E_0}]=0
 \quad\text{in }\operatorname{Pic}(E_0)\simeq
 \operatorname{Pic}(\mathrm{dP}_6).
\]
This proof uses neither global Leray
nor local integral specialization.  See
\cref{prop:local-cusp-picard-obstruction,prop:analytic-ah-boundary,%
lem:analytic-support-specialization,thm:analytic-all-picard}.

\item \Verified\ \textbf{The constructed threefold is a theorem-level
counterexample to CDP20 Theorem~2.2 as stated.}
It has
\[
H^1(X;\mathbb Z)=H^2(X;\mathbb Z)=0,\qquad
a(X)=1,\qquad f:X\to\mathbb P^1
\text{ a holomorphic algebraic reduction},
\]
while
\[
c_3(X)=2>0,\qquad X\cong_{\mathrm{diff}}S^6.
\]
These are precisely the relevant hypotheses and the opposite of the
published conclusions \(c_3(X)\le0\) and \(a(X)=0\).  This conclusion uses
the verified candidate proof, not merely the fact that defects were found in
the published proof.

\item \Verified\ \textbf{The strongest corrected CDP conclusion is now
proved, with its exact extra hypothesis.}
For any compact complex threefold in the CDP holomorphic-reduction setting,
existence of a single \(L_0\) with
\[
 H^2(X,T_X\otimes L_0)=0
\]
implies, for general \(M\),
\[
 H^i(X,T_X\otimes M)=0\quad(i\ne1),\qquad
 \chi(X,T_X\otimes M)=\frac12c_3(X)\le0.
\]
The relative condition \(R^2f_*(T_X\otimes L)=0\) for one \(L\) implies
this global hypothesis by projection formula, Serre vanishing on the base,
and Leray.  On a reduced full fibre \(W_c\), after the two endpoint
vanishings, the missing ambient-cotangent vanishing is equivalent to
injectivity of
\[
 H^0(W_c,\operatorname{tors}\Omega^1_{W_c}\otimes A_c)
 \xrightarrow{\delta_c}
 H^1(W_c,N^*_{W_c/X}\otimes A_c).
\]
This repairs the proof route without asserting the false original theorem.
Because \(c_3(X)=2\), the contrapositive of the corrected global theorem gives
\(H^2(X,T_X\otimes L_0)\ne0\) for every \(L_0\).  Independently, the candidate
has nonzero \(\ker\delta_c\) at its cusp and
\(R^2f_*(T_X\otimes L)\ne0\) for every \(L\).  See
\cref{thm:corrected-cdp-global,cor:corrected-cdp-relative,%
thm:corrected-cdp-fibrewise}.

\item \Verified\ \textbf{The Proposition~10.8 uniform non-torsion branch
is complete.}
The reconstructed direct-image calculation gives
\[
R^1f_*\mathcal O_X\simeq\mathcal O\oplus\mathcal O(-1),
\qquad \operatorname{Pic}(X)\simeq\mathbb C,
\]
with degree-one base change on all scheme-theoretic fibres.  At a finite
filling, the invariant canonical section and the direct-image generator are
not conflated: if \(k_j=m_j-1-a_j\), then
\[
 \operatorname{div}(\overline\sigma_j)=k_jS_j,
 \qquad
 f_{j*}\omega_{N_j/D_j}
 =\mathcal O_{D_j}\cdot s_j^{k_j}\vartheta'_j,
 \qquad
 s_j^{k_j}\vartheta'_j
 =\frac{e^{\varphi(s_j)}}{m_j}s_j^{-a_j}e
 \quad\text{on }\Delta_j^*.
\]
In particular \(k_2=2\), so the canonical section is not a frame there;
the invariant-module identity is the exact statement that proves the Hodge
line used below.  The finite
character calculation gives
\(\operatorname{ord}(\mathcal O_X(S_j)|_{S_j})=m_j\).  For
\(\lambda=a\gamma+bu+cw+d\delta\), membership of its Dolbeault image in
the global Picard line is exactly
\[
c\tau+d\mu-b=0.
\]
Linear independence of \(1,\tau,\mu\) makes the rank-two exceptional
fibres countable, while rank-one fibres contribute bad parameters only in
\(\mathbb R\).  The three full special fibres and the reductions \(S_j\)
each add only countably many parameters.  Thus every
\(z\in\mathbb C\setminus(\mathbb R\cup\text{a countable set})\) supplies
one \(L_z\) non-torsion on every full fibre and on both \(S_j\).
The cusp step was first repaired using intrinsic nearby cycles and therefore
does not depend on the later cellular-collapse verification.

\item \Verified\ \textbf{The equivariant cellular collapse and its exact
comparison with nearby cycles are now complete.}
The positive toric quotient is a compact $3$-manifold cobordism between two
tori, its central boundary is a strong deformation retract, and the lifted
retraction spreads by compact phases and conjugates back to the original
twist.  On the fixed fibre, the resulting map $c=r|_F$ has
\[
0\longrightarrow\operatorname{im}(\textstyle\bigwedge^qM_0-I)
\longrightarrow\textstyle\bigwedge^q\Lambda
\xrightarrow{c_*}H_q(W;\mathbb Z)\longrightarrow0.
\]
In degree two, with the primitive orientations fixed in
\cref{prop:verified-collapse-homology},
\[
[c_*]_2=
\begin{pmatrix}
1&0&0&0&0&0\\
0&0&-1&0&0&0\\
0&-1&-1&1&1&0\\
0&0&0&1&0&0
\end{pmatrix},
\]
and its kernel is
$\mathbb Z\langle\widehat w\widehat\delta,
\widehat\gamma\widehat w+\widehat u\widehat\delta\rangle$.
The neighborhood realization of the nearby-cycle unit then gives, for all
$q$,
\[
\operatorname{sp}^q=i_F^*(i_W^*)^{-1}=i_F^*r^*=c^*.
\]
Thus C1--C3 are verified integrally, including saturation and the
degree-three orientation.

\item \Verified\ \textbf{The global Mayer--Vietoris realization and its
Leray comparison are complete.}
The earlier audit conflated the two stages.  On
$X^\circ=X_1\cup X_2$, the verified radial markings give the actual maps
\[
 \alpha_q=(\pi_{1*},-\pi_{2*}).
\]
Their degree-one and degree-two matrices have saturated rank three.  On the
second attachment $X=N'_0\cup X^\circ$, primitive splittings reduce the
global maps to
\[
 D_1=\begin{pmatrix}1&0&0\\0&1&0\\-12&0&-p\end{pmatrix},
 \qquad
 D_2^{\rm red}=\begin{pmatrix}-p&0\\0&1\end{pmatrix}.
\]
For $p=-1$ they are unimodular.  Dualizing the same oriented collar maps
gives commuting obstruction squares and
\[
 d_2(12\gamma)=-p\omega,\quad
 d_2(2q)=-p\omega[\delta],\quad
 d_2(2\gamma uw)=-p\omega[\gamma\delta].
\]
Thus the alternate route independently gives the integral homology of
$S^6$; see
\cref{prop:actual-finite-band-maps,prop:cusp-mv-blocks,%
thm:mv-leray-comparison,cor:candidate-mv-sphere}.

\item \Verified\ \textbf{The Hodge diamond and Fr\"olicher sequence are
computed completely.}
Put
\[
 \Sigma=p_0+p_1+p_2,\qquad
 D=f^{-1}(\Sigma)_{\rm red}=W+S_1+S_2,\qquad
 \mathcal L=\Omega_B^1(\log\Sigma)\simeq\mathcal O_B(1).
\]
The logarithmic relative bundle
$\mathcal E=\Omega_X^1(\log D)/f^*\Omega_B^1(\log\Sigma)$ has
\[
f_*\mathcal E\simeq\mathcal O\oplus\mathcal O(-1),\qquad
R^1f_*\mathcal E\simeq\mathcal O^{\oplus2}\oplus\mathcal O(-1)^{\oplus2},
\qquad R^2f_*\mathcal E\simeq\mathcal O\oplus\mathcal O(-1).
\]
All three are locally free and commute with base change.
At a finite point write
$t=s^m$ and $\ell=dt/t=m\,ds/s$, with $(m,a)=(3,2)$ at $p_1$ and
$(m,a)=(4,1)$ at $p_2$, and let
$D_{\rm KS}=(d_{rs})$ be the marked Kodaira--Spencer matrix after removing
$ds$.  Dots in the raw period determinant mean $d/ds$, and
\[
 K=\det(\dot Z)=6\dot\mu^{\,2}-\dot\tau\dot\beta.
\]
With $\zeta_m=e^{-2\pi i/m}$, choose relative-cotangent eigenframes
$(e_0,e_1)$ with characters
$(1,\zeta_m^a)$ and $R^1\mathcal O$ eigenframes $(b_0,b_1)$ with characters
$(1,\zeta_m^{-a})$.  The connecting maps are typed as
\[
 \kappa_0:f_*\mathcal E\longrightarrow
 \mathcal L\otimes R^1f_*\mathcal O_X,\qquad
 \kappa_1:R^1f_*\mathcal E\longrightarrow
 \mathcal L\otimes R^2f_*\mathcal O_X\simeq\mathcal O.
\]
In the actual invariant source basis $(e_0,s^{m-a}e_1)$ and target basis
$(b_0,s^ab_1)\ell$, put
\[
 A=\frac{s}{m}d_{00},\quad B=\frac{s^{1-a}}m d_{01},\quad
 C=\frac{s^{1-a}}m d_{10},\quad
 D_0=\frac{s^{m+1-2a}}m d_{11}.
\]
The two connecting maps then have the exact matrices
\[
 [\kappa_0]=\begin{pmatrix}A&tB\\ C&D_0\end{pmatrix},
 \qquad
 [\kappa_1]=(-C,A,-D_0,B),
\]
where the ordered source of the row is
$(e_0b_0,s^ae_0b_1,s^{m-a}e_1b_0,e_1b_1)$ and its target generator is
$s^a(b_0\wedge b_1)\ell$.  Thus
\[
 \operatorname{div}(\det\kappa_0)=p_1+p_2,\qquad
 \operatorname{ord}_{z_1}K=2,\quad \operatorname{ord}_{z_2}K=0,\qquad
 \kappa_1\text{ is sheaf-surjective}.
\]
The residue normalizations are also exact:
\[
 \mathcal O(1)\longrightarrow
 \mathbb C_{p_0}\oplus\mathbb C_{p_1}\oplus\mathbb C_{p_2},
 \qquad
 \sigma\longmapsto(\sigma(p_0),3\sigma(p_1),4\sigma(p_2)),
\]
\[
 \operatorname{res}_{1,j}(xb_0+ys^ab_1)=m_jx[b_0],
\]
and
\[
 \operatorname{im}(\kappa_0\bmod t)
 =\ker\operatorname{res}_{1,j}=\mathbb C\,s^ab_1.
\]
Hence the two finite cokernel skyscrapers map isomorphically to
$H^1(S_1,\mathcal O)\oplus H^1(S_2,\mathcal O)$ and cancel in the absolute
residue sequence.

At the cusp, the quotient from ordinary to logarithmic relative
differentials is checked on the $xy$ and $xyz$ modules.  In the fixed global
branch ordering its two triple-point maps are
\[
 d_P(a)=(a_0-a_1,a_0-a_2,a_1-a_2),\qquad
 d_Q(a)=(a_2-a_1,a_0-a_1,a_0-a_2),
\]
Put $\eta_1=d\log x_1$ and $\eta_2=d\log x_2$.  These cusp logarithmic
forms lie in the logarithmic-form module, whereas the finite-filling
covectors denoted $\eta_j$ above lie in the filling-dual lattice.  The exact
morphism used below is the branch-difference map $\rho$ from the former
module to the branch-difference lattice; no identification of the two source
modules is made.
Both branch maps give
\[
 \rho(\eta_1)=q_1=(-1,0,1),\qquad
 \rho(\eta_2)=q_2=(0,-1,-1),
\]
which is the degree-zero branch-difference isomorphism.  With
$\alpha_i=\delta_0(q_i)$ and the standard \v Cech differential, the
total-complex sign is
\[
 \delta_1\bigl(\rho_*(\alpha\otimes\eta_i)\bigr)
 =-\alpha\smile\delta_0(q_i)=-\alpha\smile\alpha_i.
\]
Thus, for
$\theta=-\delta_1^{-1}(\alpha_1\smile\alpha_2)$, the induced row is exactly
\[
 \rho_*=(0,1,-1,0):
 \mathbb C\langle
 \alpha_1\eta_1,\alpha_1\eta_2,\alpha_2\eta_1,\alpha_2\eta_2
 \rangle\longrightarrow\mathbb C\theta.
\]
The same local modules show that $\Omega^1_{X/B}$ is torsion-free but
non-locally-free at the cusp; its torsion is supported precisely on the two
multiple fibres.

For the remaining scalar, rescale the Cayley coordinate to
$s=(z_2-\bar z_2)s_2$, so that
$s\circ g_2=-is$ and $ds/dz(z_2)=1$; write a dot for $d/ds$, and put
$\ell=dt/t=4\,ds/s$.  In the marked fibre coordinate
$\zeta=(\zeta_1,\zeta_2)$, the invariant section
$\xi=\partial/\partial\zeta_2$ is nowhere zero in the locally free relative
sheaf $T_{X/B}=\mathcal E^\vee$ (although its ambient tangent-space value can
vanish at a toric fixed point), and the unique lift of its dual unit is
$h=(a(z),1)(d\zeta_1,d\zeta_2)^{\mathsf t}$ with $a(z_2)=1/2$.
At the order-four lift the genuine invariant target frame is
\[
 \omega_2=s\,u_2(s)(\gamma\wedge u)\ell,\qquad u_2(0)=1;
\]
here $b_1$ is normalized by
$\gamma\wedge b_1=u_2(s)(\gamma\wedge u)$.
$s(\gamma\wedge u)\ell$ alone is not invariant away from $s=0$.
In this frame the signed coarse value is
\[
 \kappa_1(h\otimes\gamma)(p_2)
 =-\frac14\bigl(a(z_2)\dot\tau(z_2)+\dot\mu(z_2)\bigr).
\]
Since $\operatorname{ord}_{z_2}(\tau-i)=2$, one has
$\dot\tau(z_2)=0$ and
$K(z_2)=6\dot\mu(z_2)^2\ne0$; hence the scalar is exactly
$-\frac14\dot\mu(z_2)\ne0$.  Therefore
\[
f_*\Omega_X^1\simeq\mathcal O(-2),\quad
R^1f_*\Omega_X^1\simeq\mathcal O\oplus\mathcal O(-1)^2,\quad
R^2f_*\Omega_X^1\simeq\mathcal O\oplus\mathcal O(-1),
\]
and the exact Hodge table is
\[
\begin{array}{c|cccc}
q\backslash p&0&1&2&3\\ \hline
3&0&0&0&1\\
2&0&1&2&1\\
1&1&2&1&0\\
0&1&0&0&0.
\end{array}
\]
The two middle Fr\"olicher rows are exact copies of
$0\to\mathbb C\to\mathbb C^2\to\mathbb C\to0$.  Thus the sequence does
not degenerate at $E_1$, while
$\dim E_1=10$, $E_2=E_\infty$, and $\dim E_2=2$.  See
 \cref{prop:log-relative-direct-images,prop:log-ks-divisor,%
 prop:absolute-cotangent-direct-images,thm:h12-scalar,cor:exact-hodge-diamond}.
The internal calculation is \Verified, and a further context-isolated AI
derivation has independently recovered the invariant lattices, determinant
divisor, residue cancellation, scalar test, direct-image splittings, and the
same values $h^{1,2}=1$, $h^{1,1}=2$.  The separate H2 audit record remains
\Open\ only because independent human specialist replication is still
outstanding; that external replication boundary is not a premise of the
Hodge theorem.

\item \Verified\ \textbf{Every Bott--Chern and Aeppli dimension is fixed.}
The auxiliary spaces in McHugh's compact-threefold formulas satisfy
$k^{p,q}\le b_{p+q}$: the de Rham image of the pure, doubly closed forms
surjects onto $K^{p,q}$.  Hence they vanish in all middle degrees for this
$X$; also
$E_2^{3,1}=0$.  Thus
\[
h^{1,1}_{BC}=2,\quad h^{1,2}_{BC}=h^{2,1}_{BC}=1,\quad
h^{2,2}_{BC}=2,\quad h^{2,3}_{BC}=h^{3,2}_{BC}=1,
\]
besides the two endpoint copies of $\mathbb C$, and all other Bott--Chern
dimensions vanish.  Schweitzer duality gives the complete Aeppli diamond.
Equivalently,
\[
(h^k_{BC})=(1,0,2,2,2,2,1),\quad
(h^k_A)=(1,2,2,2,2,0,1),\quad
(\Delta^k)=(0,2,4,4,4,2,0).
\]
See \cref{thm:exact-bott-chern-aeppli}.  Coordinate-level representatives
and the product structure are certified separately in the next item.

\item \Verified\ \textbf{The Bott--Chern algebra has explicit generators
and square-zero positive ideal.}
Choose global $\bar\partial$-closed smooth representatives
$G\in\mathcal A^{0,1}(X)$ and $Q\in\mathcal A^{1,1}(X)$ of the two
candidate-specific Dolbeault classes $\mathfrak g$ and $\mathfrak q$.
Here $\mathfrak g$ is distinguished, while $\mathfrak q$ is a chosen
coherent-Leray lift, unique modulo $\operatorname{im}d_1$.
Set
\[
 \mathsf A=[\partial G]_{BC},\qquad
 \mathsf B=[\bar\partial\bar G]_{BC}.
\]
On a marked smooth fibre,
\[
 \sigma_1=6\mu\gamma+\tau u+w,\qquad
 \sigma_2=\beta\gamma+\mu u+\delta,
\]
put
\[
 y=\operatorname{Im}\tau,\qquad m=\operatorname{Im}\mu,\qquad
 r=\operatorname{Im}\beta,\qquad
 \Delta=r-\frac{6m^2}{y}<0,
\]
and define
\[
 \Gamma=\frac{i}{2\Delta}
 \left(d\bar\zeta_2-\frac{m}{y}d\bar\zeta_1\right),\qquad
 U=\frac{i}{2y\Delta}
 \left(r\,d\bar\zeta_1-6m\,d\bar\zeta_2\right).
\]
These represent the fibre images of \(\gamma\) and \(u\), and the global
classes satisfy
\[
 i_b^*\mathfrak g=[\Gamma]_{\bar\partial},\qquad
 i_b^*\mathfrak q=[Q_F]_{\bar\partial}.
\]
The Appell--Humbert matrix itself is
\[
 H_\eta=\frac1{y\Delta}
 \begin{pmatrix}r&-6m\\-6m&6y\end{pmatrix},
 \qquad H_\eta=2G_\eta^{-1},\qquad
 \det H_\eta=\frac6{y\Delta}<0.
\]
The load-bearing wedge identity is
\[
 \eta=u\wedge w+6\gamma\wedge\delta
 =u\wedge\sigma_1+6\gamma\wedge\sigma_2
 =-\sigma_1\wedge u-6\sigma_2\wedge\gamma,
\]
which gives the exact fibre representative
\[
 Q_F=-d\zeta_1\wedge U-6d\zeta_2\wedge\Gamma,
 \qquad
 \eta\wedge\eta=12\gamma\wedge u\wedge w\wedge\delta.
\]
The marked and complex orientations are opposite because the real
determinant of the marked lattice columns is
$y\Delta<0$.  Accordingly,
\[
 \int_{(F_b,{\rm or}_{\rm mark})}Q_F^2=12,
 \qquad
 \int_{(F_b,{\rm or}_{\mathbb C})}Q_F^2=-12.
\]
The fibre Dolbeault--Serre pairing is normalized by
\[
 \int_{F_b}Q_F\wedge\Gamma\wedge d\zeta_2=-1
 \quad\text{for the complex orientation},
\]
so $[\Gamma\wedge Q_F]_{\bar\partial}\ne0$.

The classes $d_1\mathfrak g$ and $\mathfrak q$ are independent in
$H^{1,1}_{\bar\partial}(X)$: the first restricts to zero on every smooth
fibre, whereas the second restricts to the nonzero class of $Q_F$.
Exactness of the first middle Fr\"olicher row therefore gives
\[
 d_1\mathfrak q=[\partial Q]_{\bar\partial}\ne0,
\]
so $\mathsf C=[\partial Q]_{BC}$ and its conjugate are nonzero.  Moreover,
putting
\[
 v=[G\wedge Q]_{\bar\partial}\in H^{1,2}_{\bar\partial}(X)
\]
gives
$i_b^*v=[\Gamma\wedge Q_F]_{\bar\partial}\ne0$.  Thus $v$ spans
$H^{1,2}_{\bar\partial}(X)$, and injectivity in the second middle row gives
\[
 d_1v=[\partial(G\wedge Q)]_{\bar\partial}\ne0.
\]
Consequently $\mathsf D=[\partial(G\wedge Q)]_{BC}$ and its conjugate are
nonzero.

The global integral normalization is equally explicit.  If the base
orientation is chosen so that
$d_2(12\gamma)=c_1(\mathcal O_B(1))$ and
$\exp(z)=e^{2\pi iz}$, the holomorphic structure
$\bar\partial_zs=-2\pi izG\otimes s$ on the smooth trivial line has
unit-metric Chern connection and curvature
\[
 \nabla_z=d-2\pi i(zG+\bar z\,\bar G),\qquad
 F_{\nabla_z}=-2\pi i
 (z\,\partial G+\bar z\,\bar\partial\bar G).
\]
With the convention
$c_{1,BC}(L,h)=[\,iF_{\nabla^h}/(2\pi)\,]_{BC}$, one obtains
\[
 \exp(12\mathfrak g)=f^*\mathcal O_B(1),\qquad
 c_{1,BC}(\exp(z\mathfrak g))
 =z\mathsf A+\bar z\,\mathsf B,
\]
and hence
\[
 c_{1,BC}(f^*\mathcal O_B(1))=12(\mathsf A+\mathsf B).
\]
With the opposite primitive orientation both occurrences of $12$ acquire
the common sign $-1$:
\[
 \exp(-12\mathfrak g)=f^*\mathcal O_B(1),\qquad
 c_{1,BC}(f^*\mathcal O_B(1))=-12(\mathsf A+\mathsf B).
\]
Thus the canonical integral Picard--Leray datum fixes the absolute factor
$12$; it does not assert that $\mathsf A,\mathsf B$ are integral
Bott--Chern classes.

Put
\[
 \mathbf1=[1]_{BC},\qquad
 \mathsf T=[\nu_X]_{BC},
 \quad \nu_X\text{ a smooth positive }(3,3)\text{-form with }
 \int_X\nu_X=1.
\]
Together with the unit, the following nine positive-degree classes form
bases of all positive-degree Bott--Chern groups, with $\mathsf T$ the
normalized top class:
\[
 \mathsf A=[\partial G]_{BC},\ \mathsf B=[\bar\partial\bar G]_{BC},\
 \mathsf C=[\partial Q]_{BC},\ \overline{\mathsf C},\
 \mathsf D=[\partial(G\wedge Q)]_{BC},\ \overline{\mathsf D},\
 \mathsf E=[\partial(Q^2)]_{BC},\ \overline{\mathsf E},\ \mathsf T.
\]
If $N$ denotes their span, then
\[
 H^{*,*}_{BC}(X)=\mathbb C\mathbf1\oplus N,
 \qquad N^2=0.
\]
For the last one-dimensional generator,
$[Q^2]_{\bar\partial}$ has nonzero fibre restriction, whereas every class
in
$\operatorname{im}(d_1:H^{1,2}_{\bar\partial}\to
H^{2,2}_{\bar\partial})$ restricts to zero.  Exactness therefore gives
\[
 d_1[Q^2]_{\bar\partial}
 =[\partial(Q^2)]_{\bar\partial}\ne0,
\]
which certifies $\mathsf E$ (and conjugation certifies
$\overline{\mathsf E}$).
Equivalently every product of two positive-degree classes is zero.  Put
$a=\partial G$, $a^\dagger=\bar\partial\bar G$, and $c=\partial Q$.
The exact primitive
\[
 \partial\bar\partial(G\wedge\bar G)=-a\wedge a^\dagger
\]
gives $\mathsf A\mathsf B=0$.  Represent
$c_{1,BC}(f^*\mathcal O_B(1))$ by the pullback Fubini--Study Chern form;
its square vanishes pointwise because $B$ is a curve.  Squaring the
factor-$12$ identity and using $\mathsf A\mathsf B=0$ gives
$\mathsf A^2+\mathsf B^2=0$.  Moreover
\[
 [(a^\dagger)^2]_{\bar\partial}=0,\qquad
 [a^2]_{\bar\partial}=d_1[G\wedge a]_{\bar\partial}=0.
\]
Injectivity of
$d_1:H^{1,2}_{\bar\partial}\to H^{2,2}_{\bar\partial}$ then gives a form
$\Phi$ with
\[
 G\wedge a=\bar\partial\Phi,\qquad
 a^2=\partial\bar\partial\Phi,
\]
so $\mathsf A^2=\mathsf B^2=0$.  The last below-top product requiring a
pairing is disposed of by
\[
 \int_Xa\wedge c\wedge G
 =-\int_XG\wedge a\wedge c
 =-\int_X\bar\partial\Phi\wedge c=0,
\]
together with the perfect Dolbeault--Serre pairing and the isomorphism
$H^{3,2}_{BC}\to H^{3,2}_{\bar\partial}$; the
Bott--Chern--Aeppli perfect pairing alone would not prove this vanishing.
Every top-degree product of two positive-degree generators contains a de
Rham exact factor, while
\[
 H^{3,3}_{BC}(X)\xrightarrow{\ \int_X\ }\mathbb C
\]
is an isomorphism, so those products vanish as well.  See
\cref{prop:bcpr-generators-are-bases,thm:bcpr-square-zero-algebra}.  This
asserts global smooth representatives, but not a preferred literal
coefficient atlas through the two finite fillings and the cusp:
$[i_b^*Q]_{\bar\partial}=[Q_F]_{\bar\partial}$ does not require
$i_b^*Q=Q_F$ as forms.  The chosen $Q$ cannot be real, although its fibre
class is real.  After fixing the distinguished $\mathfrak g$ and one
coherent-Leray lift $\mathfrak q$, changing the smooth representatives and
horizontal corrections changes $G,Q$ by Dolbeault coboundaries and changes
none of the displayed classes, pairings, or products.  This is a
candidate-specific rank-one line-bundle normalization, not a statement
about arbitrary gauge theories.
\item \Verified\ \textbf{The positive radial cusp has an exact flat
period, inverse-dual, and scalar-gap certificate.}
Put
\[
 \mu_T=\alpha_T+i m_T,\qquad h_T=\eta_T+i H_T,\qquad
 b_T=\xi_T+i B_T,\qquad c_T=\xi_T-\eta_T,
\]
\[
 L_T=T+H_T,\qquad q_T=L_T-B_T,\qquad
 D_T=L_Tq_T+6m_T^2.
\]
Set also
\[
 m=\operatorname{Im}\mu(0),\qquad
 H=\operatorname{Im}h(0),\qquad
 B=\operatorname{Im}b(0).
\]
In the ordered marked lattice and interleaved real coordinate bases,
\[
 P_T=
 \begin{pmatrix}
  6\alpha_T&\eta_T&1&0\\
  6m_T&L_T&0&0\\
  c_T&\alpha_T&0&1\\
  -q_T&m_T&0&0
 \end{pmatrix},\qquad
 \det P_T=-D_T,\qquad
 \operatorname{covol}(P_T\mathbb Z^4)=D_T>0,\qquad
 D_{\mathrm{can},T}=-\frac{D_T}{L_T}<0,
\]
and the complete inverse transpose is
\[
 P_T^{-\mathsf T}=
 \begin{pmatrix}
 0&0&1&0\\[2pt]
 \dfrac{m_T}{D_T}&\dfrac{q_T}{D_T}&
 \dfrac{-6m_T\alpha_T-q_T\eta_T}{D_T}&
 \dfrac{-m_Tc_T-q_T\alpha_T}{D_T}\\[7pt]
 0&0&0&1\\[2pt]
 -\dfrac{L_T}{D_T}&\dfrac{6m_T}{D_T}&
 \dfrac{6(L_T\alpha_T-m_T\eta_T)}{D_T}&
 \dfrac{L_Tc_T-6m_T\alpha_T}{D_T}
 \end{pmatrix}.
\]
Consequently
\((P_T\mathbb Z^4)^*=P_T^{-\mathsf T}\mathbb Z^4\).  More explicitly,
for \(k=(r,s,n_1,n_2)^{\mathsf t}\), define
\[
 R_T(k)=r-6\alpha_Tn_1-c_Tn_2,\qquad
 S_T(k)=s-\eta_Tn_1-\alpha_Tn_2.
\]
The corresponding dual vector is
\[
 \left(n_1,
 \frac{m_TR_T(k)+q_TS_T(k)}{D_T},n_2,
 \frac{-L_TR_T(k)+6m_TS_T(k)}{D_T}\right),
\]
and the scalar eigenvalue is \(4\pi^2\) times its squared norm.
For every sufficiently large \(T\),
\[
 \lambda_1(F_T,g_T)
 =4\pi^2\min\left\{
 \frac{L_T^2+m_T^2}{D_T^2},
 \frac{q_T^2+36m_T^2}{D_T^2}\right\}
 =\frac{4\pi^2}{T^2}\left(
 1+\frac{-2H+2\min\{B,0\}}T+O(T^{-2})\right).
\]
This is the gap of the explicitly declared Euclidean quotient metric.  Its
unit-volume rescaling is a different operator family and obeys
\[
 \widehat g_T=D_T^{-1/2}g_T,\qquad
 \lambda_1(F_T,\widehat g_T)
 =\frac{4\pi^2}{T}\left(
 1+\frac{-H-B/2+2\min\{B,0\}}T+O(T^{-2})\right)
 \longrightarrow0.
\]
The degeneration has an exact fixed-energy density, not merely a lowest-mode
estimate.  Counting Fourier multiplicities by
\[
 N_T(E)=\#\{k\in\mathbb Z^4:
 4\pi^2\lVert P_T^{-\mathsf T}k\rVert^2\le E\},
 \qquad R_E^2=\frac{E}{4\pi^2},
\]
one has, for every fixed \(E\ge0\),
\[
 N_T(E)=\pi D_T\sum_{n\in\mathbb Z^2}
 (R_E^2-\lVert n\rVert^2)_+ +O_E(T+1),
 \qquad (x)_+=\max\{x,0\}.
\]
The sum is finite.  In particular, below the first transverse threshold,
\[
 0<E<4\pi^2
 \quad\Longrightarrow\quad
 N_T(E)=\frac{E}{4\pi}D_T+O_E(T+1)
       =\frac{E}{4\pi}T^2+O_E(T).
\]
The proof writes the admissible \((r,s)\)-indices for each fixed transverse
\(n\in\mathbb Z^2\) as the lattice points in a translated ellipse of area
\(\pi D_T(R_E^2-\lVert n\rVert^2)\) and perimeter \(O_E(T+1)\), then applies
the elementary planar area--perimeter lattice bound.  Thus the remainder is
uniform in the otherwise unrestricted real parts of the periods.  See
\cref{cusp-spec-counting}.

Normalized Haar measure in the angular-coordinate direct integral is a
unitary measure convention, not this metric rescaling.  See
 \cref{cusp-spec-exact-dual-lattice,cusp-spec-shortest-modes,%
 cusp-spec-unit-volume,cusp-spec-counting}.

\item \Verified\ \textbf{The native
\(\mathbb T^2\times\mathbb R^2\) limit and the ancillary covered
\(\mathbb R^4\) limit have different exact constructions.}
The real projection \(\rho(x_1,y_1,x_2,y_2)=(y_1,y_2)\) gives
\[
 0\longrightarrow\mathbb T^2_{\mathrm{re}}
 \longrightarrow F_T
 \longrightarrow Q_T:=\mathbb R^2/Y_T\mathbb Z^2
 \longrightarrow0,
 \qquad
 Y_T=\begin{pmatrix}6m_T&L_T\\-q_T&m_T\end{pmatrix},
 \qquad \det Y_T=D_T.
\]
The quotient systole is at least \(T-C\), whereas the kernel retains its
two unit periods; this is the native pointed
\(\mathbb T^2\times\mathbb R^2\) degeneration.  Adding
\[
 \mathsf S_M=\operatorname{diag}(1,1,M,M),\qquad
 t_c=v^M
\]
produces a degree-\(M^2\) relative fibre cover over \(v\ne0\), because
\[
 \mathsf S_M^{-1}M_0^M\mathsf S_M=M_0.
\]
The cover character indexed by \(k\in\mathbb Z^4\) has momentum
\[
 P_T^{-\mathsf T}\mathsf S_M^{-1}k,
\]
and it is pulled back from the original fibre exactly when
\(k_3,k_4\) are divisible by \(M\).  The covering radius of this momentum
 lattice is \(O(T^{-1}+M^{-1})\); hence \(T,M\to\infty\) fills all four
 momentum directions and gives eventual exact Euclidean balls.  This argument
 by itself proves only a pointed local \(\mathbb R^4\) limit over the
 punctured base; the modified-action completion across the singular fibre is
 the distinct local Kummer construction recorded below, and neither argument
 alone gives a global cover of the completed threefold or an S12/S24-type
 globalization.

There is, however, an exact four-copy regular-locus construction.  The
monodromy-invariant closure of the compact cusp lattice is
\[
 L_\Lambda=\mathbb Z\langle\widehat u,\widehat w,\widehat\delta\rangle,
\qquad
 \Lambda/L_\Lambda\cong\mathbb Z\langle\widehat\gamma\rangle.
\]
The induced real circle quotient has exact circumference
\[
 \ell_T=\frac{D_T}{\sqrt{L_T^2+m_T^2}}
     =T+H-B+O(T^{-1})\longrightarrow\infty.
\]
The fourfold fibre product of \(X^\circ=f^{-1}(B^\circ)\) therefore has a
real fibrewise-group \(\mathbb T^4\) quotient over \(B^\circ\) whose
synchronized radial fibres converge pointed-smoothly to \(\mathbb R^4\).
Two copies are insufficient for this monodromy-invariant-kernel mechanism;
four are minimal, and no finite cover of the base repairs the native two-copy
quotient.  The two finite fillings admit explicit simultaneous-root
normalizations over Seifert targets, with degrees \(27\) and \(64\) on the
reduced central fibres.  The toric cusp does not: two paths converging to the
same triple point retain different character values \(3/4\) and \(1/2\), so
the literal quotient and its fourfold product have no continuous cusp
 extension.  The Cartesian source is not \(S^{24}\):
 \(X^4\simeq(S^6)^4\), while that fixed-target fourfold-product map does not
 factor through the topological smash sphere
\(X^{\wedge4}\cong S^{24}\); Borel--Serre also rules out almost-complex
structures on \(S^{12}\) and \(S^{24}\).  Thus this exact regular-locus
construction answers the globalization question as far as this mechanism can
go and proves its compact-extension obstruction.  See
 \cref{cusp-spec-native-quotient,cusp-spec-globalization-obstruction,%
 cusp-spec-cover-momentum-theorem}.

\item \Verified\ \textbf{The degree-\(N^2\) regular-fibre cover has an exact
completed local toric Kummer model, with direct and relative degrees kept
distinct.}
Write
\[
 Z(s)=sB_0+C(t),\qquad
 B_0=\begin{pmatrix}0&1\\-1&0\end{pmatrix},\qquad
 t=e^{2\pi i s},
\]
and let the target lattice action on the toric variety \(Y=Y_\Sigma\) be
\[
 \Psi_{\bar\lambda}(x,t)
 =\bigl(e^{2\pi iC(t)\bar\lambda}
        t^{B_0\bar\lambda}x,t\bigr),
 \qquad N_0=Y/\bar\Lambda.
\]
After the base change \(t=v^N\), the correct source factor and action are
\[
 C_N(v)=\frac1N C(v^N),\qquad
 \Psi^{(N)}_{\bar\lambda}(\xi,v)
 =\bigl(e^{2\pi iC_N(v)\bar\lambda}
        v^{B_0\bar\lambda}\xi,v\bigr).
\]
They extend freely and properly discontinuously to a copy \(Y^{(N)}\) of
the same toric variety.  Thus
\[
 N_0^{\langle N\rangle}=Y^{(N)}/\bar\Lambda
 \longrightarrow\Delta_v
\]
is a proper holomorphic map from a smooth threefold, with reduced irreducible
nonnormal central fibre \(W_N\), whose normalization is \(\mathrm{dP}_6\).  The toric
power map
\[
 P_N(\xi_1,\xi_2,v)=(\xi_1^N,\xi_2^N,v^N)
\]
is finite flat of degree \(N^3\), is equivariant for the two displayed
actions, and descends to
\[
 \kappa_N:N_0^{\langle N\rangle}\longrightarrow N_0,
 \qquad \deg\kappa_N=N^3,\qquad
 \kappa_N^*W=NW_N,\qquad
 \operatorname{Ram}(\kappa_N)=(N-1)W_N.
\]
If \(N_0[N]=N_0\times_{\Delta_t}\Delta_v\), then \(N_0[N]\) is normal and
\[
 N_0^{\langle N\rangle}\xrightarrow{Q_N}N_0[N]\longrightarrow N_0
\]
is the unique relative factorization.  The generic degree of \(Q_N\) is
\(N^2\); for \(N>1\) it is quasi-etale but not flat, with non-etale and
nonflat loci exactly over the double and triple strata of the base-changed
central fibre.  It is not a normalization.  On reduced central fibres the
map has generic degree \(N^2\), and on normalizations it is the toric map
\[
 [N]:\mathrm{dP}_6\longrightarrow\mathrm{dP}_6,
\]
of degree \(N^2\), ramified with index \(N\) along all six toric boundary
curves and respecting the paired conductor curves and the two triple-point
classes.  Over \(v\ne0\), it is exactly the lattice cover
\[
 \mathbb Z\langle\widehat\gamma,\widehat u,
 N\widehat w,N\widehat\delta\rangle\subset\Lambda,
 \qquad \operatorname{Deck}(Q_N)\cong(\mathbb Z/N\mathbb Z)^2,
 \qquad S_N^{-1}M_0^NS_N=M_0,
\]
where \(S_N=\operatorname{diag}(1,1,N,N)\).  The factor
\(C(v^N)/N\) is essential: reusing the target action would require the
generally false equality
\(e^{2\pi iC(v)\bar\lambda N}=e^{2\pi iC(v^N)\bar\lambda}\).
This is a morphism of completed \emph{local} cusp families.  For \(N>1\)
the cover lattice is not preserved by the full monodromy, so it gives
neither a cover over all of \(B\) nor a product- or smash-sphere map.  See
\cref{cusp-cover-extension-theorem,cusp-cover-naive-action-failure,%
cusp-cover-global-scope}.

\item \Verified\ \textbf{The maps \(q_2\), \(q_{1,\mathrm{cusp}}\), and
\(q_1\) give exact cone-smash carriers with sources \(S^{12}\) and
\(S^{24}\), and their intrinsic metrics, ambient metrics, homology, and
limitations are all explicit.}
Let
\[
 K_\Lambda=\mathbb Z\langle\widehat w,\widehat\delta\rangle,
 \qquad \overline\Lambda=\Lambda/K_\Lambda,
 \qquad \overline T^2=(\overline\Lambda\otimes\mathbb R)/\overline\Lambda.
\]
Fix \(C_{\mathrm{sys}}>0\) so that, for all sufficiently large \(T\),
\[
 Y_T=T\begin{pmatrix}0&1\\-1&0\end{pmatrix}+E_T,
 \qquad \lVert E_T\rVert_{\mathrm{op}}\le C_{\mathrm{sys}}.
\]
Since
\[
 \operatorname{im}(M_0-I)=K_\Lambda=\ker(M_0-I),
\]
the punctured cusp has a fibrewise group epimorphism
\[
 \overline\pi:N_0\setminus W\longrightarrow D_0^*\times\overline T^2
\]
with connected two-torus kernel and trivial target local system.  On the
radial fibre, \(\overline T^2\) is isometric to
\(Q_T=\mathbb R^2/Y_T\mathbb Z^2\) and has systole at least
\(T-C_{\mathrm{sys}}\).

Put
\[
 \mathcal A=\operatorname{Susp}(S^1\times\overline T^2)
 \cong\operatorname{Susp}(T^3).
\]
For \(t(f(x))=re^{i\phi}\), the continuous surjective quotient
\[
 q_2:(X,x_0)\longrightarrow(\mathcal A,v_0),\qquad
 q_2(x)=[r,e^{i\phi},\overline\pi_f(x)]
\]
on \(N_0\setminus W\), where
\(\overline\pi_f=\operatorname{pr}_{\overline T^2}\circ\overline\pi\), sends
\(W\) to \(v_0\) and
\(X\setminus\operatorname{int}N_0\) to the other suspension vertex
\(v_1\).  Its smash square
\[
 q_2^{\wedge2}:X\wedge X\cong S^{12}\longrightarrow
 \mathcal A\wedge\mathcal A
\]
restricts on synchronized radial fibres to
\[
 \pi_T\times\pi_T:F_T^2\longrightarrow Q_T^2,
 \qquad (Q_T^2,0)\longrightarrow(\mathbb R^4,0),
\]
where the convergence is pointed local-isometric, pointed smooth, and
pointed Gromov--Hausdorff, and every fixed radius ball is eventually exactly
Euclidean.  The nonzero reduced integral homology groups are
\[
 \widetilde H_k(\mathcal A)=
 \begin{cases}\mathbb Z^3,&k=2,3,\\ \mathbb Z,&k=4,\end{cases}
\quad
 \widetilde H_k(\mathcal A^{\wedge2})=
 \begin{cases}
 \mathbb Z^9,&k=4,\\ \mathbb Z^{18},&k=5,\\
 \mathbb Z^{15},&k=6,\\ \mathbb Z^6,&k=7,\\ \mathbb Z,&k=8.
 \end{cases}
\]
More precisely,
\[
 \mathcal A^{\wedge2}\simeq
 \bigvee^9S^4\vee\bigvee^{18}S^5\vee\bigvee^{15}S^6
 \vee\bigvee^6S^7\vee S^8.
\]

The cusp-supported circle version uses
\[
 \mathcal B=\operatorname{Susp}(S^1\times S^1)=\operatorname{Susp}(T^2)
\]
and the continuous surjective based quotient
\[
 q_{1,\mathrm{cusp}}:(X,x_0)\longrightarrow(\mathcal B,b_0),
 \qquad
 q_{1,\mathrm{cusp}}(x)=[r,e^{i\phi},\chi_\gamma(x)],
 \qquad q_{1,\mathrm{cusp}}(W)=b_0,
\]
with the exterior collapsed to \(b_1\).  Its fourth smash power has source
\[
 q_{1,\mathrm{cusp}}^{\wedge4}:X^{\wedge4}\cong S^{24}
 \longrightarrow\mathcal B^{\wedge4}
\]
and restricts to
\(\chi_{\gamma,T}^{\times4}:F_T^4\to C_T^4\).  Here
\[
 \widetilde H_k(\mathcal B)=
 \begin{cases}\mathbb Z^2,&k=2,\\ \mathbb Z,&k=3,\end{cases}
\quad
 \widetilde H_k(\mathcal B^{\wedge4})=
 \begin{cases}
 \mathbb Z^{16},&k=8,\\ \mathbb Z^{32},&k=9,\\
 \mathbb Z^{24},&k=10,\\ \mathbb Z^8,&k=11,\\ \mathbb Z,&k=12.
 \end{cases}
\]

For the full-base version, glue the regular circle target and the two finite
Seifert targets to obtain the compact three-manifold with boundary
\[
 \mathcal M=(P\times S^1)\cup C_1\cup C_2,
 \qquad \partial\mathcal M=\partial D_0\times S^1,
\]
and cone its boundary:
\[
 \mathcal C=\mathcal M\cup_{\partial\mathcal M}
 \operatorname{Cone}(\partial\mathcal M).
\]
The map \(q_1:X\to\mathcal C\) is \((f,\theta_\gamma)\) over \(P\),
equivalently \(\chi_\gamma|_{f^{-1}(P)}\), is the
finite-filling map \(N_j\to C_j\) for \(j=1,2\), and on the cusp is
\[
 q_1(x)=[r,(e^{i\phi},\theta_\gamma(x))],\qquad q_1(W)=c_0.
\]
It is continuous, surjective, and based.  Its fourth smash power is
\[
 q_1^{\wedge4}:X^{\wedge4}\cong S^{24}\longrightarrow\mathcal C^{\wedge4},
\]
and its synchronized radial restriction is again
\(\chi_{\gamma,T}^{\times4}:F_T^4\to C_T^4\), with
\[
 \operatorname{sys}(C_T^4)=\ell_T=T+H-B+O(T^{-1}),
 \qquad (C_T^4,0)\longrightarrow(\mathbb R^4,0).
\]
The target homology is
\[
 \widetilde H_k(\mathcal C)=
 \begin{cases}\mathbb Z,&k=2,3,\end{cases}
\quad
 \widetilde H_k(\mathcal C^{\wedge4})=
 \begin{cases}
 \mathbb Z,&k=8,12,\\ \mathbb Z^4,&k=9,11,\\
 \mathbb Z^6,&k=10.
 \end{cases}
\]
In fact,
\[
 \mathcal C^{\wedge4}\simeq
 S^8\vee\bigvee^4S^9\vee\bigvee^6S^{10}
 \vee\bigvee^4S^{11}\vee S^{12}.
\]
All omitted degrees in these three homology displays are zero.
The conical-target maps \(q_2^{\wedge2}\) and \(q_1^{\wedge4}\) induce zero on
all positive-degree integral homology and cohomology and have zero ordinary
Hurewicz image; neither has a continuous homotopy section.  Their unstable
homotopy classes remain \Open, because those facts do not decide nullity.  The
class of the cusp-supported conical-target map
\(q_{1,\mathrm{cusp}}^{\wedge4}:S^{24}\to\mathcal B^{\wedge4}\) has likewise
not been computed.

The two-copy and four-copy exponents are sharp for the unscaled quotient
norm induced by the marked Euclidean fibre metric.  If a primitive lattice
\(N_{\mathrm{quot}}\subset\Lambda^{\oplus r}\) defines a connected quotient
over the punctured cusp whose radial systole tends to infinity, then every
bounded-length class must vanish, so
\[
 K_\Lambda^{\oplus r}\subset N_{\mathrm{quot}},\qquad
 \operatorname{rank}(\Lambda^{\oplus r}/N_{\mathrm{quot}})\le2r.
\]
Thus a rank-four quotient needs at least two copies, and equality is attained
by \((\Lambda/K_\Lambda)^{\oplus2}\).  Over the full regular base, including
after a finite base cover, monodromy invariance and primitivity additionally
force \(\widehat u\in N_{\mathrm{quot}}\): indeed
\(A_1\widehat w=\widehat u-\widehat w\), while a finite-index subgroup contains
some power \(G_*^m\) satisfying
\(G_*^m\widehat w=\widehat w+m(\widehat u-\widehat\delta)\).
Consequently
\[
 \operatorname{rank}(\Lambda^{\oplus r}/N_{\mathrm{quot}})\le r,
\]
so four copies are necessary and sufficient there.  This proves the exact
minimality of the \(S^{12}\) cusp-local and \(S^{24}\) full-base smash
sources within this native connected torus-quotient class; it asserts no
minimality among arbitrary continuous quotient maps or after changing to a
\(T\)-dependent rescaling.  See \cref{cusp-cone-optimal-exponents}.

There are also sphere targets, not merely conical targets.  Let
\[
 K^{(1)}=((\mathbb R/\mathbb Z)\times\{0\})
 \cup(\{0\}\times(\mathbb R/\mathbb Z)),\qquad
 c:\overline T^2\twoheadrightarrow\overline T^2/K^{(1)}\cong S^2,
\]
where \(c\) collapses the coordinate one-skeleton.  Forgetting the cusp base
angle in the reduced suspensions gives continuous surjections
\[
 \widehat p_2:\mathcal A\twoheadrightarrow S^3,
 \quad[r,e^{i\phi},u]\longmapsto[r,c(u)],
 \qquad
 \widehat p_1:\mathcal B\twoheadrightarrow S^2,
 \quad[r,e^{i\phi},\theta]\longmapsto[r,\theta].
\]
Thus \(\widehat q_3=\widehat p_2\circ q_2\) and
\(\widehat q_\gamma=\widehat p_1\circ q_{1,\mathrm{cusp}}\) induce, after the
chosen based diffeomorphism \(X\cong_{\mathrm{diff}}S^6\), quotient maps
\[
 \widehat q_3^{\wedge2}:S^{12}\twoheadrightarrow S^6,
 \qquad
 \widehat q_\gamma^{\wedge4}:S^{24}\twoheadrightarrow S^8.
\]
Both full sphere carriers are based null-homotopic, by two distinct exact
calculations.  First let \(f:S^6\to S^2\) be any based map.  The factorization
\[
 S^{12}=S^6\wedge S^6
 \xrightarrow{\ f\wedge\operatorname{id}_{S^6}=\Sigma^6f\ }
 S^2\wedge S^6=S^8
 \xrightarrow{\ \operatorname{id}_{S^2}\wedge f\ }
 S^2\wedge S^2=S^4
\]
and
\(\pi_{12}(S^8)\cong\pi_4^{\mathrm S}=0\) show that \(f\wedge f\) is
based null-homotopic.  Associativity gives
\(f^{\wedge r}=(f\wedge f)\wedge f^{\wedge(r-2)}\), so
\[
 [f^{\wedge r}]=0\in\pi_{6r}(S^{2r})
 \qquad(r\ge2).
\]
Taking \(f=\widehat q_\gamma\) proves, in particular,
\[
 [\widehat q_\gamma^{\wedge2}]=0\in\pi_{12}(S^4),
 \qquad
 [\widehat q_\gamma^{\wedge4}]=0\in\pi_{24}(S^8).
\]
Thus the full surjective quotient \(S^{24}\twoheadrightarrow S^8\) is based
null-homotopic.

For the other carrier, let \(g:S^6\to S^3\) be any based map.  Its exact
smash factorization is
\[
 g\wedge g=(\Sigma^3g)\circ(\Sigma^6g):
 S^{12}\longrightarrow S^9\longrightarrow S^6.
\]
Writing \(\operatorname{st}\) for stabilization, the low-stem data are
\[
 \pi_6(S^3)\cong\mathbb Z/12,
 \quad
 \pi_3^{\mathrm S}=\mathbb Z/24\langle\nu\rangle,
 \quad
 \pi_6^{\mathrm S}=\mathbb Z/2\langle\nu^2\rangle,
\]
\[
 \pi_{12}(S^6)=\mathbb Z/2\langle\nu_6\circ\nu_9\rangle,
 \qquad
 \operatorname{st}:\pi_{12}(S^6)
 \xrightarrow{\ \cong\ }\pi_6^{\mathrm S}.
\]
The full groups are tabulated in \cite[p.~339]{Hatcher2002}; the signed smash
formula and the unstable stabilization calculation are
\cite[Propositions~3.1 and~5.11]{Toda1962}.
Because \(12[g]=0\), its stable image is
\(\operatorname{st}([g])=2a\nu\) for some \(a\in\mathbb Z/12\).  Stable
smash multiplication gives
\[
 \operatorname{st}([g\wedge g])
 =\operatorname{st}([g])^2
 =4a^2\nu^2=0\in\pi_6^{\mathrm S}.
\]
The displayed stabilization isomorphism forces
\([g\wedge g]=0\in\pi_{12}(S^6)\).
Associativity again shows that \(g^{\wedge r}\) is based null-homotopic for
every \(r\ge2\).  Taking \(g=\widehat q_3\) therefore proves
\[
 [\widehat q_3^{\wedge2}]=0\in\pi_{12}(S^6),
 \qquad
 [\widehat q_\gamma^{\wedge4}]=0\in\pi_{24}(S^8).
\]
Both maps retain their point-set surjectivity and quotient-map status.  The
one-copy class \([\widehat q_3]\in\pi_6(S^3)\) is now zero by
Theorem~\ref{h24:thm:cusp-g-zero}, while
\([\widehat q_\gamma]\in\pi_6(S^2)\) remains open. Neither full-map
null theorem settles any of the conical-target classes named above; see
\cref{cusp-cone-S12-sphere-map-null,cusp-cone-S24-sphere-map-null,%
cusp-cone-sphere-map-homotopy-status,cusp-cone-target-homotopy-types}.
Give \(S_T^2=\overline T^2/K^{(1)}\) the quotient length metric transported
from \(Q_T\), and give both smash products the quotient length metrics from
the Riemannian \(\ell^2\)-product metrics on \(Q_T^2\) and \(C_T^4\).
These moving-slice metrics are assigned intrinsically and are not restrictions
of one fixed ambient sphere metric.  The targets contain moving topologically
embedded four-spheres
\[
 \mathbb S_{12,T}^4=S_T^2\wedge S_T^2\subset S^6,
 \qquad
 \mathbb S_{24,T}^4=C_T^{\wedge4}\subset S^8,
\]
and the synchronized radial restrictions are exactly
\[
 F_T^2\xrightarrow{\pi_T^{\times2}}Q_T^2
 \xrightarrow{c_T^{\times2}}(S_T^2)^2
 \longrightarrow\mathbb S_{12,T}^4,
 \qquad
 F_T^4\xrightarrow{\chi_{\gamma,T}^{\times4}}C_T^4
 \longrightarrow\mathbb S_{24,T}^4,
\]
where the last arrows collapse the corresponding fat wedges.  The full maps
from \(F_T^2\) and \(F_T^4\) retain the nontrivial kernels of their first
arrows; at the quotient stage, precisely the displayed skeleton/fat-wedge
loci are collapsed.  At the
midpoint basepoints, every radius-\(R\) ball is exactly a Euclidean four-ball
once \(T-C_{\mathrm{sys}}>4R\), respectively \(\ell_T>4R\).  Hence both
moving four-spheres converge to \(\mathbb R^4\) in pointed local-isometric and
pointed Gromov--Hausdorff topology; on those eventual Euclidean basepoint
neighborhoods the convergence is also smooth.  Each moving four-sphere is
null-homotopic in its ambient sphere by sliding its suspension coordinate to
an endpoint, so its fundamental class dies in the target.  See
\cref{cusp-cone-angle-forgotten-sphere-carriers,%
cusp-cone-moving-S4-null-homotopy}.
Removing the smash basepoint gives exact homeomorphisms
\[
 (Q_T\setminus K_T^{(1)})^2
 \cong\mathbb S_{12,T}^4\setminus\{*\}
 \cong\mathbb R^4\hookrightarrow S^6,
\qquad
 (C_T\setminus\{0\})^4
 \cong\mathbb S_{24,T}^4\setminus\{*\}
 \cong\mathbb R^4\hookrightarrow S^8.
\]
These are nonclosed topological open cells, not globally Euclidean-isometric
subspaces for a fixed compatible compact-target metric; see
\cref{cusp-cone-literal-R4-open-cells}.
Both sphere-to-sphere quotient maps induce zero on every positive-degree
integral homology group, and both full maps are based null-homotopic:
\[
 [\widehat q_3^{\wedge2}]=0\in\pi_{12}(S^6),
 \qquad
 [\widehat q_\gamma^{\wedge4}]=0\in\pi_{24}(S^8).
\]
The one-copy map $\widehat q_3$ is now null-homotopic by
Theorem~\ref{h24:thm:cusp-g-zero}; the one-copy class of
$\widehat q_\gamma$ remains open. These full-map statements are
distinct from the two explicit null-homotopies of the moving four-sphere
inclusions.

For completeness, the radial target-slice embeddings used below are typed as
follows.  Put \(r_T=e^{-2\pi T}\) and
\[
 \vartheta_T:\mathbb R/\mathbb Z\xrightarrow{\ \cong\ }C_T,
 \qquad[a]\longmapsto[\ell_Ta].
\]
Then
\[
 j_T:Q_T^2\hookrightarrow\mathcal A^{\wedge2},\qquad
 j_T(y_1,y_2)=
 \bigl[[r_T,1,\iota_T^{-1}(y_1)],
       [r_T,1,\iota_T^{-1}(y_2)]\bigr]_{\wedge},
\]
\[
 k_T^{\mathcal B}:C_T^4\hookrightarrow\mathcal B^{\wedge4},\qquad
 k_T^{\mathcal B}(c_1,\ldots,c_4)=
 \bigl[[r_T,1,\vartheta_T^{-1}(c_1)],\ldots,
       [r_T,1,\vartheta_T^{-1}(c_4)]\bigr]_{\wedge},
\]
and, in the cusp-cone piece of \(\mathcal C\),
\[
 k_T^{\mathcal C}:C_T^4\hookrightarrow\mathcal C^{\wedge4},\qquad
 k_T^{\mathcal C}(c_1,\ldots,c_4)=
 \bigwedge_{j=1}^4[r_T,(1,\vartheta_T^{-1}(c_j))].
\]
These are the embeddings of
\cref{cusp-cone-target-slice-embeddings}.

The intrinsic decompactification can be made into an actual boundary family.
Fix \(T_0\) sufficiently large, put \(r_0=e^{-2\pi T_0}\), and define the
compact radial carriers
\[
 \mathcal D_{12}
 =\frac{[0,r_0]\times(\overline T^2)^2}
 {\{0\}\times(\overline T^2)^2\sim d_{12}},
 \qquad
 \mathcal D_{24}
 =\frac{[0,r_0]\times(S^1)^4}
 {\{0\}\times(S^1)^4\sim d_{24}}.
\]
Both are homeomorphic to \(\operatorname{Cone}(\mathbb T^4)\).  At either
cone vertex \(d_\nu\),
\[
 H_k(\mathcal D_\nu,\mathcal D_\nu\setminus\{d_\nu\};\mathbb Z)
 \cong
 \begin{cases}
  \mathbb Z^4,&k=2,\\
  \mathbb Z^6,&k=3,\\
  \mathbb Z^4,&k=4,\\
  \mathbb Z,&k=5,\\
  0,&\text{otherwise},
 \end{cases}
\]
so \(d_\nu\) is not a topological-manifold point, with or without boundary,
of any dimension.  Set
\[
 A_T^{(12)}=\operatorname{diag}(Y_T,Y_T),\qquad
 A_T^{(24)}=\ell_T I_4,\qquad s_0=T_0^{-1},
\]
and, for \(\nu\in\{12,24\}\), quotient
\([0,s_0]\times\mathbb R^4\) fibrewise by
\[
 (0,x)\sim_\nu(0,y)\iff x=y,\qquad
 (s,x)\sim_\nu(s,y)\iff
 x-y\in A_{1/s}^{(\nu)}\mathbb Z^4\quad(s>0).
\]
The resulting \(\mathfrak E_\nu\) is a Hausdorff, second-countable topological
five-manifold with boundary, with
\[
 p_{12}^{-1}(s)\cong Q_{1/s}^2,\qquad
 p_{24}^{-1}(s)\cong C_{1/s}^4\quad(s>0),\qquad
 p_\nu^{-1}(0)\cong\mathbb R^4.
\]
Its other boundary component, over \(s=s_0\), is a compact four-torus.
Every fixed Euclidean ball embeds isometrically in all sufficiently small
positive fibres.  There is a continuous surjective comparison map
\(\beta_\nu:\mathfrak E_\nu\to\mathcal D_\nu\) which is a homeomorphism away
from the zero fibre, identifies each positive fibre with its radial torus
slice, and collapses the entire zero-boundary component
\(p_\nu^{-1}(0)\cong\mathbb R^4\) to the vertex of the
compact
\(\mathcal D_\nu\cong\operatorname{Cone}(\mathbb T^4)\).  Thus the Euclidean
space is an exact decompactification boundary, not the whole compact source or
target and not a continuous quotient target of a compact source.  The map
\(\beta_\nu\) is neither proper nor a quotient map, so this is not a
topological blow-down or a quotient identification
\(\mathfrak E_\nu/\mathbb R^4\cong\mathcal D_\nu\).  Explicitly, for
\(a=(1/2,1/2,1/2,1/2)\) and \(s_j\downarrow0\), the points
\[
 z_j=[s_j,A_{1/s_j}^{(\nu)}a]
\]
form a closed subset of \(\mathfrak E_\nu\).  Indeed, in the \(S^{12}\) case
\[
 \inf_{n\in\mathbb Z^4}\lVert A_{1/s_j}^{(12)}(a+n)\rVert
 \ge(1/s_j-C_{\mathrm{sys}})\operatorname{dist}(a,\mathbb Z^4)\to\infty,
\]
and in the \(S^{24}\) case the same infimum equals
\(\ell_{1/s_j}\operatorname{dist}(a,\mathbb Z^4)\to\infty\); hence the
sequence has no accumulation in a positive fibre and eventually misses every
half-space chart about every zero-fibre point.  On the other hand,
\(\beta_\nu(z_j)=[e^{-2\pi/s_j},a]\to d_\nu\).  Hence the image set is not
closed although its inverse image is closed, which proves that
\(\beta_\nu\) is not quotient;
see \cref{cusp-cone-exact-decompactification-space,%
cusp-cone-sharp-global-R4-scope}.

Unrescaled ambient metrics impose a different conclusion.  Put
\[
 Z_{12}:=\mathcal A^{\wedge2},\qquad
 Z_{24}^{\mathcal B}:=\mathcal B^{\wedge4},\qquad
 Z_{24}^{\mathcal C}:=\mathcal C^{\wedge4}.
\]
For every fixed compatible metric on these compact targets,
\[
 \operatorname{diam}_{\mathrm{amb}}j_T(Q_T^2)\to0,
 \qquad
 \operatorname{diam}_{\mathrm{amb}}k_T^{\mathcal B}(C_T^4)\to0,
 \qquad
 \operatorname{diam}_{\mathrm{amb}}k_T^{\mathcal C}(C_T^4)\to0,
\]
whereas intrinsically
\[
 \operatorname{diam}(Q_T^2)\ge\tfrac12(T-C_{\mathrm{sys}})\to\infty,
 \qquad \operatorname{diam}(C_T^4)\ge\tfrac12\ell_T\to\infty.
\]
There are fixed cusp-adapted compatible length metrics with regular-stratum
metric tensors \(g_{Z_{12}}\) and \(g_{Z_{24}^{\mathcal Q}}\),
\(\mathcal Q\in\{\mathcal B,\mathcal C\}\), for which the exact rescaled first
fundamental forms are
\[
 \left(\frac{T}{r_T}\right)^2j_T^*g_{Z_{12}}
 =g_{Q_T}\oplus g_{Q_T},
 \qquad
 \left(\frac{\ell_T}{r_T}\right)^2
 (k_T^{\mathcal Q})^*g_{Z_{24}^{\mathcal Q}}
 =g_{C_T}^{\oplus4},
 \quad\mathcal Q\in\{\mathcal B,\mathcal C\}.
\]
With their induced Riemannian path metrics (not the restricted ambient
distance functions), the rescaled slices converge to \(\mathbb R^4\), but the corresponding full
ambient pointed blow-ups converge on their eventual regular-stratum balls to
\(\mathbb R^8\) and \(\mathbb R^{12}\), respectively.  In every rescaling
that unrolls a slice to \(\mathbb R^4\), the cone/fat-wedge basepoint recedes
to infinite scaled distance; vertex-based conic tangents have dimensions
\(8\) and \(12\) for exact conic target metrics, and unit-volume normalization
gives a flat four-torus, not
\(\mathbb R^4\).  See
\cref{sphere-metric-uniform-collapse,sphere-metric-fixed-realization,%
sphere-metric-scaling-trichotomy,sphere-metric-unit-volume,%
sphere-metric-tangent-scope}.

Thus coning changes the codomain and does not contradict the proved
nonextension of the literal character map to the fixed target
\(B\times S^1\).  Nor do Cartesian powers give spheres:
\[
 X^2\simeq S^6\times S^6\not\simeq S^{12},\qquad
 X^4\simeq(S^6)^4\not\simeq S^{24};
\]
only the based smash powers are spheres.  There is no full-base two-copy
four-torus quotient, no fixed \(T^4\) target retaining the four character
classes, and no fixed \(S^4\) target extending a degree-one collapse of the
radial quotient.  No nontrivial finite base cover preserves
\(K_\Lambda\), and the local unspooling cover is not a nontrivial connected
unbranched cover of the simply connected compact sources.  Finally,
\(S^{12}\) and \(S^{24}\) admit no almost-complex structure, so the smash
sources are not compact complex spheres.  See
\cref{cusp-cone-S12-route,cusp-cone-S24-local-route,cusp-cone-S24-route,%
cusp-cone-metric-scope,cusp-cone-obstructions,%
cusp-cone-explicit-carriers,cusp-cone-target-homotopy-types}.
Every construction in this sphere-route item uses the native quotients of the
original, unscaled cusp fibres; none uses the ancillary finite-cover
unspooling from the preceding item.

\item \Verified\ \textbf{The repaired compact-link construction and its
finite Gaussian quantization have exact actions, spectra, and a resolved
zero-gap diagonal sequence.}
Keep the cover index and the mesh index independent:
\[
 M_{\mathrm{cov}}\in\mathbb Z_{\ge1},\qquad
 S_{M_{\mathrm{cov}}}=\operatorname{diag}(1,1,M_{\mathrm{cov}},M_{\mathrm{cov}}),
 \qquad P_{T,M_{\mathrm{cov}}}=P_TS_{M_{\mathrm{cov}}},
 \qquad N_{\mathrm{mesh}}\in\mathbb Z_{\ge2}.
\]
Then
\[
 \widetilde F_{T,M_{\mathrm{cov}}}
 =\mathbb R^4/P_{T,M_{\mathrm{cov}}}\mathbb Z^4
 \longrightarrow F_T
\]
has degree \(M_{\mathrm{cov}}^2\), Gram matrix and covolume
\[
 G_{T,M_{\mathrm{cov}}}
 =P_{T,M_{\mathrm{cov}}}^{\mathsf T}P_{T,M_{\mathrm{cov}}},
 \qquad
 \mathcal D_{T,M_{\mathrm{cov}}}
 =\sqrt{\det G_{T,M_{\mathrm{cov}}}}
 =M_{\mathrm{cov}}^2D_T.
\]
Put
\[
 V_{N_{\mathrm{mesh}}}
 =(\mathbb Z/N_{\mathrm{mesh}}\mathbb Z)^4,
 \qquad
 E_{N_{\mathrm{mesh}}}
 =V_{N_{\mathrm{mesh}}}\times\{1,2,3,4\}.
\]
Transport of the standard
\(N_{\mathrm{mesh}}\)-cubulation by \([y]\mapsto[P_{T,M_{\mathrm{cov}}}y]\)
is an honest cubulation of the whole cover torus; its vertex set is exactly
the \(N_{\mathrm{mesh}}\)-torsion subgroup, via
\[
 [n]\longmapsto
 \left[P_{T,M_{\mathrm{cov}}}n/N_{\mathrm{mesh}}\right].
\]
Parallel transport along its oriented edges gives the typed gauge-equivariant
holonomy map to \(\mathrm{SU}(2)^{E_{N_{\mathrm{mesh}}}}\).

Write \(M=M_{\mathrm{cov}}\) and \(N=N_{\mathrm{mesh}}\) only within the
following compact-link formulas.  Normalize
\[
 \langle X,Y\rangle_{\mathfrak{su}(2)}
 =-\tfrac12\operatorname{tr}(XY),
 \qquad H_{\mathrm{link}}=\operatorname{diag}(i,-i),
 \qquad \lVert H_{\mathrm{link}}\rVert^2=1,
 \qquad \exp(2\pi H_{\mathrm{link}})=I_2.
\]
The periodic
inverse-transport links
\[
 \begin{aligned}
 U_1(n)&=\exp(-2\pi n_2H_{\mathrm{link}}/N^2),\\
 U_2(n)&=
 \begin{cases}I_2,&0\le n_2\le N-2,\\
 \exp(2\pi n_1H_{\mathrm{link}}/N),&n_2=N-1,\end{cases}\\
 U_3(n)&=U_4(n)=I_2.
 \end{aligned}
\]
have, including every wrap plaquette,
\[
 U_{12}(n)=\exp(2\pi H_{\mathrm{link}}/N^2),\qquad U_{ij}(n)=I_2
 \quad((i,j)\ne(1,2),\ i<j),
\]
The principal-logarithm action and the
explicit continuum-matched \((1,2)\)-Wilson sector are exactly
\[
 S^{\log}_{T,M,N}(U)=\frac{2\pi^2M^2}{g_{\mathrm{YM}}^2D_T},
\]
\[
 S^{W,12}_{T,M,N}(U)
 =\frac{M^2N^4}{2g_{\mathrm{YM}}^2D_T}
   \left(2-2\cos\frac{2\pi}{N^2}\right)
 =\frac{M^2}{g_{\mathrm{YM}}^2D_T}
  \left(2\pi^2-\frac{2\pi^4}{3N^4}
  +\frac{4\pi^6}{45N^8}+O(N^{-12})\right).
\]
The inverse-Gram two-form minor producing these formulas is exactly
\[
 (G_{T,M_{\mathrm{cov}}}^{-1})_{11}
 (G_{T,M_{\mathrm{cov}}}^{-1})_{22}
 -(G_{T,M_{\mathrm{cov}}}^{-1})_{12}^2=D_T^{-2}.
\]

For \(0\le p\le4\), define the real cochain spaces
\[
 C_N^p
 =\operatorname{Map}\!\left(
 V_{N_{\mathrm{mesh}}},
 \bigwedge\nolimits^p(\mathbb R^4)^*\otimes\mathfrak{su}(2)\right).
\]
For \(\vartheta^j_{T,M}=P_{T,M}^{-\mathsf T}\mathbf e_j\), let
\(\mathsf T_j a(n)=a(n+\mathbf e_j)\) and set
\[
 D^+_{N,j}=N(\mathsf T_j-I),
 \qquad D^-_{N,j}a(n)=N(a(n)-a(n-\mathbf e_j)),
\]
\[
 d_Na=\sum_{j=1}^4
 \varepsilon(\vartheta^j_{T,M})D^+_{N,j}a,
 \qquad
 d_N^*=-\sum_{j=1}^4
 \iota_{(\vartheta^j_{T,M})^\sharp}D^-_{N,j},
\]
and
\[
 \Delta_N^{(1)}=d_Nd_N^*+d_N^*d_N,
 \qquad
 \mathscr C_{T,M,N}
 =\ker(d_N^*:C_N^1\to C_N^0)
 \cap(\ker\Delta_N^{(1)})^\perp.
\]
This last orthogonal complement removes the zero Fourier harmonic cochains.
On this real reduced Coulomb space put
\[
 \mathcal L_{T,M,N}=g_{\mathrm{YM}}^{-2}d_N^*d_N,
 \qquad \mathfrak d_N(k)_j=N(e^{2\pi ik_j/N}-1).
\]
For every \(k\ne0\), its eigenvalue and complex transverse multiplicity are
\[
 \lambda_{T,M,N}(k)=\frac1{g_{\mathrm{YM}}^2}
 \left\lVert P_T^{-\mathsf T}S_{M_{\mathrm{cov}}}^{-1}
 \mathfrak d_N(k)\right\rVert^2,
 \qquad 3\dim\mathfrak{su}(2)=9.
\]
In particular,
\[
 \lambda_{T,M,N}(\mathbf e_1)
 =\frac{4N^2\sin^2(\pi/N)}{g_{\mathrm{YM}}^2}
  \frac{L_T^2+m_T^2}{D_T^2},
\]
independently of \(M\).  The Friedrichs realization on
\(L^2(\mathscr C_{T,M,N},da)\) of
\[
 \widehat H^{\mathrm{Gauss}}_{T,M,N}
 =-\tfrac12\Delta_{\mathscr C_{T,M,N}}
  +\tfrac12\langle a,\mathcal L_{T,M,N}a\rangle
\]
has a unique Gaussian ground state and exact gap
\[
 \operatorname{gap}(\widehat H^{\mathrm{Gauss}}_{T,M,N})
 =\sqrt{\lambda_{\min}(\mathcal L_{T,M,N})}
 \le\frac{2N\sin(\pi/N)}{g_{\mathrm{YM}}}
       \frac{\sqrt{L_T^2+m_T^2}}{D_T}.
\]
For
\[
 T_j=j^2,\qquad M_j=j,\qquad N_j=j^3,
\]
one has simultaneously
\[
 \operatorname{sys}(\widetilde F_{T_j,M_j})\to\infty,
 \qquad
 \max_r\frac{\lVert P_{T_j,M_j}\mathbf e_r\rVert}{N_j}\to0,
 \qquad
 \operatorname{gap}(\widehat H^{\mathrm{Gauss}}_{T_j,M_j,N_j})\to0,
 \qquad
 S^{\log}_{T_j,M_j,N_j}(U)\to0.
\]
The exponential map \(a\mapsto(\exp(a_e/N))_e\) is only a local
diffeomorphism at zero, has derivative \(N^{-1}\operatorname{id}\), is not
globally injective, and, if the compact-link action satisfies
\(DS_W(\mathbf1)=0\), transports its Hessian with the exact factor
\[
 D^2(S_W\circ\operatorname{Exp}_N)_0
 =N^{-2}D^2(S_W)_{\mathbf1}.
\]
Consequently the displayed Gaussian Hilbert space is
not the Haar \(L^2\)-space on compact links or its nonlinear gauge-invariant
subspace; no unitary identification, nonlinear compact-link Hamiltonian,
Gauss-law spectral problem, transfer matrix, renormalized continuum limit,
or Osterwalder--Schrader reconstruction has been constructed.  This theorem
therefore remains strictly on the auxiliary finite Gaussian side of the
physical Yang--Mills boundary.  See
 \cref{cusp-link-cubical-realization,cusp-link-holonomy-map,%
 cusp-link-flux-theorem,cusp-link-finite-spectrum,cusp-link-gaussian-theorem,%
 cusp-link-exponential-local-only,cusp-link-physical-hilbert-space,%
 cusp-link-dirac-not-l2}.

\item \Verified\ \textbf{The exceptional first-coordinate elliptic quotient
is finite and does not approach the cusp.}
For the marked complex torus
\[
 \mathcal A_z=\mathbb C^2/
 \mathbb Z\langle(6\mu,\beta),(\tau,\mu),(1,0),(0,1)\rangle,
\]
the first coordinate descends onto
\(E_\tau=\mathbb C/(\mathbb Z\tau+\mathbb Z)\) exactly on
\[
 6\mu=a\tau+b,\qquad(a,b)\in\mathbb Z^2.
\]
The kernel is connected and equals
\(\mathbb C/(\mathbb Z+\mathbb Z(a\mu-\beta))\), with
\(\operatorname{Im}(a\mu-\beta)=-D_{\mathrm{can}}>0\).  Under the two
generators the labels transform by
\[
 (a,b)\xmapsto{g_0}(a,a+b),\qquad
 (a,b)\xmapsto{g_2}(-b,6+a).
\]
The marked union is therefore monodromy invariant; its image in the completed
orbifold base is finite and avoids a punctured cusp neighborhood.  Thus this
specific elliptic quotient cannot realize the cusp gap collapse.  It remains
logically possible for a fibre to have a different elliptic quotient or an
integral magnetic class not induced by this projection.  See
\cref{cusp-spec-first-coordinate-descent,%
cusp-spec-elliptic-descent-finiteness,%
cusp-spec-elliptic-descent-scope}.

\item \Verified\ \textbf{A single specified auxiliary operator has
zero essential spectral threshold.}
Fourier expansion identifies
\[
 \mathscr H_{\mathrm{aux}}
 =L^2\bigl([T_0,\infty),dT;
 L^2_0(\mathbb R^4/\mathbb Z^4,dy)\bigr)
 \cong\bigoplus_{k\in\mathbb Z^4\setminus\{0\}}
 L^2([T_0,\infty),dT).
\]
With
\[
 \lambda_k(T)=4\pi^2\lvert P_T^{-\mathsf T}k\rvert^2,\qquad
 (\mathscr L_{\mathrm{aux}}u)_k(T)=\lambda_k(T)u_k(T),
\]
the exact domain is
\[
 D(\mathscr L_{\mathrm{aux}})=\left\{u:
 \sum_{k\ne0}\int_{T_0}^{\infty}
 \lambda_k(T)^2\lvert u_k(T)\rvert^2\,dT<\infty\right\}.
\]
The operator is nonnegative and self-adjoint, has zero kernel, and the
\(k=(1,0,0,0)\) singular sequence supported on \([n,n+1]\) proves
\[
 0\in\sigma_{\mathrm{ess}}(\mathscr L_{\mathrm{aux}}).
\]
The word ``gapless'' here belongs to this auxiliary Hilbert space and
operator; it is not an invariant of the complex threefold alone.  See
\cref{cusp-interpretation-direct-integral-gapless,%
cusp-interpretation-fixed-complex-counterexample}.

\item \Verified\ \textbf{The continuum and finite classical gauge
fluctuation operators have exact spectra and cusp gaps.}
For a compact Lie group with \(d_G=\dim_{\mathbb R}\mathfrak g\), the
product-connection Hessian on the reduced Coulomb space is
\[
 \mathcal L_T=g_{\mathrm{YM}}^{-2}\mathrm d^*\mathrm d,\qquad
 \lambda_T^{\mathrm{YM}}(k)
 =\frac{4\pi^2}{g_{\mathrm{YM}}^2}
  \lvert P_T^{-\mathsf T}k\rvert^2
\]
on every \(k\ne0\) transverse Fourier fibre, with complex multiplicity
\(3d_G\).  Hence
\[
 \gamma_T=\frac{4\pi^2}{g_{\mathrm{YM}}^2}
 \min\left\{\frac{L_T^2+m_T^2}{D_T^2},
                  \frac{q_T^2+36m_T^2}{D_T^2}\right\}
 =\frac{4\pi^2}{g_{\mathrm{YM}}^2T^2}
 \left(1+\frac{-2H+2\min\{B,0\}}T+O(T^{-2})\right).
\]
For the finite cubical complex,
\[
 \mathfrak d_N(k)_j=N(e^{2\pi i k_j/N}-1),\qquad
 \zeta_{N,T}(k)=P_T^{-\mathsf T}\mathfrak d_N(k),\qquad
 \mathcal L_{N,T}=g_{\mathrm{YM}}^{-2}\mathrm d_N^*\mathrm d_N,
\]
and the reduced \(k\ne0\) eigenvalue is exactly
\[
 \lambda_{N,T}^{\mathrm{YM}}(k)
 =g_{\mathrm{YM}}^{-2}\lvert P_T^{-\mathsf T}\mathfrak d_N(k)\rvert^2
\]
with the same complex multiplicity.  Writing
\(\varrho_N^{\mathrm{spec}}=4N^2\sin^2(\pi/N)\), there is a \(T_1\), independent of
\(N\ge2\), such that for every \(T\ge T_1\) its gap is
\[
 \gamma_{N,T}=\frac{\varrho_N^{\mathrm{spec}}}{g_{\mathrm{YM}}^2}
 \min\left\{\frac{L_T^2+m_T^2}{D_T^2},
                  \frac{q_T^2+36m_T^2}{D_T^2}\right\}
 =\frac{\varrho_N^{\mathrm{spec}}}{g_{\mathrm{YM}}^2T^2}
 \left(1+\frac{-2H+2\min\{B,0\}}T+O(T^{-2})\right)
\]
uniformly for \(N\ge2\).  Thus each fixed finite model is positively
gapped, but \(\gamma_{N(T),T}\to0\) for every \(N(T)\ge2\); resolving
the varying continuum geometry is the stronger condition
\(N(T)/T\to\infty\).  The finite operator is a linearized
Lie-algebra-valued cochain model, not a nonlinear compact-link Wilson
action.  See
 \cref{cusp-gauge-finite-spectrum,cusp-gauge-finite-gap-limits,%
 cusp-gauge-continuum-symbol,cusp-gauge-continuum-gap}.

\item \Verified\ \textbf{Every fixed nondegenerate integral flux on the
marked four-torus has an exact magnetic-oscillator spectrum whose first
spacing tends to zero along the cusp.}
In angular coordinates \(u\in\mathbb R^4/\mathbb Z^4\), let
\[
 G_T=P_T^{\mathsf T}P_T,
 \qquad \sqrt{\det G_T}=D_T,
 \qquad
 \mathsf E=\sum_{j<k}E_{jk}\,\mathrm du_j\wedge\mathrm du_k
 \in\bigwedge^2(\mathbb Z^4)^*.
\]
The factor of automorphy and connection
\[
 J_n^{\mathsf E}(u)
 =\exp\!\left(-2\pi i\sum_{j<k}E_{jk}n_ju_k\right),
 \qquad
 \nabla^{\mathsf E}
 =\mathrm d+2\pi i\sum_{j<k}E_{jk}u_j\,\mathrm du_k
\]
define a Hermitian line bundle \(\mathscr L_{\mathsf E,T}\) with curvature
\(F_{\nabla^{\mathsf E}}=2\pi i\mathsf E\).  The nonnegative closed form
\[
 q_{\mathsf E,T}[f]=\int_{[0,1]^4}
 (G_T^{-1})^{jk}\overline{\nabla_j^{\mathsf E}f}
 \nabla_k^{\mathsf E}f\,D_T\,\mathrm d^4u
\]
has form domain \(H^1(F_T;\mathscr L_{\mathsf E,T})\); its Friedrichs
operator \(\Delta_{\mathsf E,T}=(\nabla^{\mathsf E})^*\nabla^{\mathsf E}\)
has domain \(H^2(F_T;\mathscr L_{\mathsf E,T})\) and compact resolvent.
For nondegenerate \(\mathsf E\), define \(b_1,b_2>0\) by
\[
 \operatorname{Spec}\!\left(
 G_T^{-1/2}(2\pi\mathsf E)G_T^{-1/2}\right)
 =\{ib_1,-ib_1,ib_2,-ib_2\}.
\]
The ladder commutators give the complete set of spectral values
\[
 \operatorname{Spec}(\Delta_{\mathsf E,T})
 =\{(2n_1+1)b_1+(2n_2+1)b_2:
 n_1,n_2\in\mathbb Z_{\ge0}\},
 \qquad
 \operatorname{gap}(\Delta_{\mathsf E,T})=2\min\{b_1,b_2\}.
\]
Writing
\[
 d=|\operatorname{Pf}(\mathsf E)|
 =|E_{12}E_{34}-E_{13}E_{24}+E_{14}E_{23}|>0,
\]
the determinant identity is
\[
 b_1b_2=\frac{4\pi^2d}{D_T},
\]
and therefore
\[
 \operatorname{gap}(\Delta_{\mathsf E,T})
 \le4\pi\sqrt{\frac d{D_T}}
 =\frac{4\pi\sqrt d}{T}
 \left(1+\frac{B-2H}{2T}+O(T^{-2})\right)\longrightarrow0.
\]
This is a fibrewise scalar magnetic Laplacian on the marked cusp tori.  It
does not assert that \([\mathsf E]\) is fixed by full monodromy or identify
this level spacing with a physical Yang--Mills particle mass.  See
\cref{cusp-magnetic-oscillator-theorem,cusp-magnetic-frequency-product}.

\item \Verified\ \textbf{A reducible non-flat \(SU(2)\) critical
connection has exact action and an exact negative Jacobi eigenvalue.}
The rank-two sum
\(\mathscr E_T=\mathscr L_T\oplus\mathscr L_T^{-1}\) has parallel trivial
determinant, \(c_1=c_2=0\), and an explicitly nullhomotopic
\(SU(2)\) clutching map.  Thus its principal \(SU(2)\)-bundle is
smoothly trivial.  In the lifted line-bundle splitting
\(\mathscr L_T\oplus\mathscr L_T^{-1}\), the direct-sum connection and its
curvature are
\[
 H_{\mathrm{mag}}=\begin{pmatrix}i&0\\0&-i\end{pmatrix},\qquad
 \nabla^T=\mathrm d+2\pi u_1\,\mathrm du_2\,H_{\mathrm{mag}},\qquad
 F_{\nabla^T}=\frac{2\pi}{D_T}H_{\mathrm{mag}}\,
 \mathrm dy_1\wedge\mathrm dy_2,
\]
where the Lie-algebra norm is the same one fixed above,
\(\langle X,Y\rangle=-\tfrac12\operatorname{tr}(XY)\), so that
\(\lVert H_{\mathrm{mag}}\rVert^2=1\).
In any chosen global \(SU(2)\)-trivialization of the smoothly trivial
principal bundle this matrix is replaced by its corresponding gauge
transform.  One has \(\mathrm d_{\nabla^T}^*F_{\nabla^T}=0\) and
\[
 \operatorname{YM}_T(\nabla)
 =\frac1{2g_{\mathrm{YM}}^2}
  \int_{F_T}|F_\nabla|^2\,\mathrm d\operatorname{vol}_{F_T},
\]
so
\[
 \operatorname{YM}_T(\nabla^T)
 =\frac{2\pi^2}{g_{\mathrm{YM}}^2D_T}
 =\frac{2\pi^2}{g_{\mathrm{YM}}^2T^2}
 \left(1+\frac{B-2H}{T}+O(T^{-2})\right).
\]
Let \((e_\mu)_{\mu=1}^4\) be a constant orthonormal frame for the descended
flat metric.  The background-gauge Jacobi expression
\[
 (\mathcal J_Ta)_\nu
 =-\sum_{\mu=1}^4\nabla_{e_\mu}\nabla_{e_\mu}a_\nu
  -2\sum_{\mu=1}^4[F_{\nu\mu},a_\mu]
\]
has an explicit smooth real unit eigenvector satisfying
\[
 \mathrm d_{\nabla^T}^*a=0,\qquad
 \mathcal J_Ta=-\frac{4\pi}{D_T}a,\qquad
 \operatorname{YM}_T''|_{\nabla^T}(a,a)
 =-\frac{4\pi}{g_{\mathrm{YM}}^2D_T}.
\]
This proves a classical saddle direction at every finite \(T\), not a
spectral assertion for a quantum Hamiltonian.  See
\cref{cusp-magnetic-su2-triviality,cusp-magnetic-ym-critical-point,%
cusp-magnetic-negative-mode}.

\item \Verified\ \textbf{The native full-flux background extends root by
root from split \(SU(2)\) to split-Cartan \(SU(r)\), with exact topology,
action, spectrum, index, and wall nullities.}
For every nonzero trace-zero \(\mathbf m\in\mathbb Z^r\), put
\[
 \mathscr E_{\mathbf m,T}
 =\bigoplus_i\mathscr L_{\eta,T}^{\otimes m_i},
 \qquad s_2(\mathbf m)=\sum_i m_i^2,
 \qquad k_{\mathbf m}=6s_2(\mathbf m).
\]
Then
\[
 c_2(\mathscr E_{\mathbf m,T})
 =-\frac{s_2(\mathbf m)}2[\eta_T]^2,
 \quad
 \langle c_2,[F_T]_{\mathrm{mark}}\rangle=-6s_2(\mathbf m),
 \quad
 \langle c_2,[F_T]_{x\text{-}\mathrm{complex}}\rangle=k_{\mathbf m},
\]
and the corresponding reducible Yang--Mills connection satisfies
\[
 \operatorname{YM}_T
 =\frac{2\pi^2k_{\mathbf m}}{g_{\mathrm{YM}}^2}
  +\frac{\pi^2s_2(\mathbf m)S_T^2}{g_{\mathrm{YM}}^2D_T}.
\]
For each positive root \(i<j\), the charge is
\(q_{ij}=|m_i-m_j|\); the complete negative normalized Coulomb bands are
\[
 q_{ij}\bigl((2n+1)\mathcal B_{T,\min}-\mathcal B_{T,\max}\bigr),
 \qquad 0\le n<N_-(T),
\]
each contributing real dimension \(12q_{ij}^2\), with coincident bands
summed.  Hence
\[
 \operatorname{ind}_{\mathbb R}\mathcal J_{\mathbf m,T}
 =12rN_-(T)s_2(\mathbf m)
 =36r s_2(\mathbf m)
\]
on the deep cusp tail.  At \(S_T=0\) the negative spectrum is empty but the
root zero modes double, so no strict-stability conclusion is made.  These are
classical fibrewise theorems, not an interacting quantum spectrum; moreover
their nonzero \(c_2\) obstructs extension over every proposed sphere carrier
with vanishing fourth integral cohomology.  See
\cref{cusp-native-full-flux-suN-topology,%
cusp-native-full-flux-suN-action-theorem,%
cusp-native-full-flux-suN-Morse-theorem,%
cusp-native-full-flux-suN-nullity-theorem,%
cusp-native-full-flux-suN-stability-boundary}.

\item \Verified\ \textbf{The Koide cone is isomorphic only to an
ambient period-coefficient discriminant, not to the regular cusp family.}
For
\[
 Q_K(r)=\frac23(r_1+r_2+r_3)^2-(r_1^2+r_2^2+r_3^2),\qquad
 \begin{pmatrix}h\\x\\y\end{pmatrix}
 =\begin{pmatrix}
 1/\sqrt3&1/\sqrt3&1/\sqrt3\\
 1/\sqrt2&0&-1/\sqrt2\\
 1/\sqrt6&-2/\sqrt6&1/\sqrt6
 \end{pmatrix}
 \begin{pmatrix}r_1\\r_2\\r_3\end{pmatrix},
\]
the coordinate change is orthogonal and
\[
 Q_K(r)=h^2-x^2-y^2.
\]
The real-linear map
\[
 (h,x,y)\longmapsto(t,m,\varrho)
 =\left(h+y,\frac{x}{\sqrt6},h-y\right)
\]
is an ambient coefficient-space isomorphism.  For the canonical convention
\[
 Z=\begin{pmatrix}6\mu&\tau\\\beta&\mu\end{pmatrix},\qquad
 P=\begin{pmatrix}
 6\operatorname{Re}\mu&\operatorname{Re}\tau&1&0\\
 6m&t&0&0\\
 \operatorname{Re}\beta&\operatorname{Re}\mu&0&1\\
 \varrho&m&0&0
 \end{pmatrix},
\]
one has, independently of the real parts,
\[
 \det P=t\varrho-6m^2
 =\det\begin{pmatrix}t&\sqrt6\,m\\\sqrt6\,m&\varrho\end{pmatrix}
 =Q_K(r).
\]
For \(r\in\mathbb R_{\ge0}^3\setminus\{0\}\), the equation
\(Q_K(r)=0\) is equivalent to the displayed symmetric matrix being
positive semidefinite of rank one.  It also forces \(\det P=0\), so
the columns of \(P\) do not form a full lattice and define no regular
compact complex two-torus.  Actual radial cusp fibres instead satisfy
\[
 \det P_T=-D_T<0,\qquad
 D_T=T^2+(2H-B)T+O(1),\qquad
 \det P_T\sim-T^2.
\]
See
\cref{koide-s6-thm-ambient-discriminant,koide-s6-cor-ambient-only,%
koide-s6-cor-actual-cusp}.

\item \Verified\ \textbf{The positive defect-one Koide shell has an exact
ideal-and-unit parametrization and an unconditional power-saving counting
law.}
Define
\[
 \Delta(x,y,z)=4(xy+xz+yz)-(x^2+y^2+z^2),
\]
and put
\[
 \mathcal R=\{(x,y,z)\in\mathbb Z_{>0}^3:
 x\le y\le z,\ \Delta(x,y,z)=-1,\ z\ge4(x+y)\},
\]
\[
 R_{\mathrm K}(Y)
 =\#\{(x,y,z)\in\mathcal R:z\le Y\}.
\]
The terminal Vieta wall is exact: below it,
the other root \(4(x+y)-z\) is positive and smaller than \(y\).  On the wall,
the coordinates
\[
n=z-4(x+y),\qquad q=y-x,\qquad a=x+y-4n
\]
give a bijection with the integral chamber
\[
a^2-3q^2-18n^2=-2,\qquad
a>0, q,n\ge0, a+4n>q,\qquad z=4a+17n,
\]
whose inverse is
\[
 (x,y,z)=
 \left(\frac{a+4n-q}{2},\frac{a+4n+q}{2},4a+17n\right).
\]
For
\[
 F_{\mathrm K}(a,q,n)=a^2-3q^2-18n^2,\qquad
 \mathscr X=\{F_{\mathrm K}=-2\},
\]
direct
Hensel counting and the good-prime formula give
\[
\sigma_2=\frac32,\qquad \sigma_3=2,\qquad
 \sigma_p=\frac{1-p^{-2}}
 {1-\left(\frac3p\right)p^{-1}}\qquad(p\ge5),
\]
and hence
\[
 \mathfrak S(\mathscr X)
 =\frac92\frac{\log(2+\sqrt3)}{\sqrt3\,\zeta(2)}.
\]
The exact Poincar\'e-residue integral over the reduced chamber is
\[
 I_\infty=\frac1{2\sqrt3}\log\!\left(\frac3{2\sqrt2}\right).
\]
The Gorodnik--Oh--Shah sector theorem applies to the genuine norm
\(\max\{|4a+17n|,|q|,|a|\}\).  An explicit element of
\(\operatorname{SO}(F_{\mathrm K},\mathbb Z_2)\) with spinor norm \(-17\) exchanges the
two Borovoi sign orbits and proves \(c_{\mathscr X}=1\).  Calibrating both
counting theorems against the same Euclidean balls therefore gives
\[
 C_0:=\frac9{2\pi^2}\log(2+\sqrt3)
 \log\!\left(\frac3{2\sqrt2}\right)
\]
and
\[
 \boxed{R_{\mathrm K}(Y)\sim C_0Y.}
\]
For \(n\geq1\), the exact ideal-and-unit parametrization puts
\[
\begin{aligned}
 m_n&=18n^2-2=2(3n-1)(3n+1),\\
 \mathcal P(m_n)&=\bigl\{(C,D,E)\in\mathbb Z_{>0}^2\times\mathbb Z:\\
 &\hspace{5em}C^2D=m_n,\quad 0\leq E<D,\quad
 E^2\equiv3\pmod D\bigr\}.
\end{aligned}
\]
The map
\[
 (C,D,E)\longmapsto
 C\bigl(D\mathbb Z+(E+\sqrt3)\mathbb Z\bigr)
\]
is a bijection to norm-\(m_n\) ideals of \(\mathbb Z[\sqrt3]\).  Each ideal
has exactly one generator \(a_0+q_0\sqrt3\) with \(a_0>2q_0\geq0\), and its
positive solutions are
\[
 a_k+q_k\sqrt3=(a_0+q_0\sqrt3)(2+\sqrt3)^k,\qquad k\geq0.
\]
The inverse descent is
\((a,q)\mapsto(2a-3q,2q-a)\), iterated while \(a\leq2q\), followed by the
unique HNF recovery.  Hence multiplicity is one and
\[
 \#\mathcal P(m_n)=\rho(m_n)
 =\sum_{d\mid m_n}\chi_{12}(d)
 =\rho(3n-1)\rho(3n+1).
\]
The \(n=0\) face is not absorbed into this positive-norm ideal family: it is
exactly
\[
 a_k+q_k\sqrt3=(1+\sqrt3)(2+\sqrt3)^k,\qquad k\geq1,
\]
with the original strict wall, and contributes the explicit floor count
\(R_0(Y)=O(\log(2+Y))\) in \cref{koide-param-eq:R0-exact}.
For every real \(Y\geq0\), the unchanged weak inequality
\(4a+17n\leq Y\) is retained in the exact floor formula of
\cref{koide-param-thm:exact-floor}; its discrepancy splits into the radial
shifted divisor term and the exact phase discrepancies.  If
\(S_\ell(m)\) denotes the norm-\(m\) coefficient of the unitary Hecke
character defined in \cref{koide-param-eq:Hecke-character}, then the angular
data obey the exact shifted identity
\[
 S_\ell(18n^2-2)
 =(-1)^\ell S_\ell(3n-1)S_\ell(3n+1),
\]
with the shifted factorization proved in
\cref{koide-param-eq:shifted-Hecke-factorization}; the endpoint half-weights
and all five sharp walls are retained in
\cref{koide-param-eq:Fejer-exact}.
A direct proof of effective well-roundedness for the unchanged sharp family
allows Benoist--Oh, Theorem~13.2~\cite[pp.~1931--1932]{BenoistOh_2012},
to be applied orbit by orbit and strengthens
the count to
\[
 \boxed{R_{\mathrm K}(Y)=C_0Y+O(Y^{1-\eta})}
 \qquad\text{for some fixed, presently unspecified }\eta>0.
\]
The primary-source comparison gives no exact-domain square-root theorem:
Liu--Sarnak's \(39/64\) exponent is for a smooth standard-radial count on
the exact congruence lattice, and the sharp standard-radial results of Good
and Chatzakos--Petridis have the \(2/3\) frontier after small-spectrum terms.
Neither transfers to the five-wall height \(4a+17n\); see
\cref{koide-lit-determination}.
The exact finite census gives
\(R_{\mathrm K}(10^8)=3{,}536{,}002\), with all rows independently replayed;
the largest \(\lvert R_{\mathrm K}(Y)-C_0Y\rvert/\sqrt Y\) for
\(10^4\leq Y\leq10^8\) is
\(0.1714001928088145597\ldots\) at \(Y=867{,}240\).

There is also an exact hyperbolic bridge.  For
\[
 B=(18,-6)_{\mathbb Q},\qquad
 \mathcal O=\mathbb Z1\oplus
 \mathbb Z\frac{I+J}{2}\oplus
 \mathbb Z\frac{I-J}{2}\oplus
 \mathbb Z\frac{IJ}{2},
\]
the map
\[
 (a,q,n)\longmapsto2+\frac{qI+aJ+nIJ}{2}
\]
is a bijection from the parity-restricted ternary shell to the
reduced-trace-\(4\), reduced-norm-\(1\) locus.  The order has reduced
discriminant \(108\), index \(18\) in a maximal overorder, and is not an
Eichler order; \(B\) is ramified exactly at \(2,3\).  Each image is primitive
hyperbolic with centralizer
\(\{\pm\gamma^k:k\in\mathbb Z\}\), eigenvalues \(2\pm\sqrt3\), and length
\(2\log(2+\sqrt3)\).  The shell is a finite union of conjugacy classes, but
its canonical height \(4a+17n\) differs from the exact radial Huber height
\(7a+15\sqrt2\,n\), since
\[
 X_i(\gamma)^2=1+\frac{(7a+15\sqrt2\,n)^2}{48}.
\]
This theorem is independent of the existence of the proposed complex
structure on \(S^6\).  The stronger square-root discrepancy
\(R_{\mathrm K}(Y)-C_0Y=O_\varepsilon(Y^{1/2+\varepsilon})\) is an explicitly \Open
problem, not a hypothesis used in the asymptotic; see
\cref{koide-count-thm-coordinate-bijection,koide-count-thm-local-densities,%
koide-count-thm-archimedean-factor,koide-count-thm-linear-asymptotic,%
koide-param-thm:integral-parametrization,koide-param-thm:power-saving,%
koide-count-sec-census-1e8,%
koide-quat-exact-reformulation-of-the-koide-chamber-count,%
koide-lit-determination,%
koide-count-open-square-root-discrepancy}.

\item \Verified\ \textbf{The earlier Erd\H{o}s--Straus divisor selector and
the \(S^6\) terminal presentation are the same marked integral matrix slice.}
For \(p\equiv1\pmod{12}\) prime, define
\[
 L=\frac{p-1}{4},\qquad 0\le j<L,\qquad
 R_j=4j+3,\qquad a_j=\frac{p+R_j}{4}.
\]
For a selector state
\[
 u\in\mathbb Z_{>0},\qquad
 u\mid a_j^2,\qquad R_j\mid4u+p^\varepsilon,\qquad
 \varepsilon\in\{0,1\},
\]
set
\[
 T=\frac{4u+p^\varepsilon}{R_j},\qquad
 \ell_0=\frac{T+1}{4},\qquad
 \ell_1=u-jT,\qquad \ell_2=1.
\]
Then \(T\equiv3\pmod4\),
\[
 12\ell_0-4\ell_1-3\ell_2=p^\varepsilon,
\]
and
\[
 \boxed{
 \begin{pmatrix}1&-j\\0&1\end{pmatrix}
 \begin{pmatrix}R_j&-u\\4&-T\end{pmatrix}
 =\begin{pmatrix}3&-\ell_1\\4&1-4\ell_0\end{pmatrix}.}
\]
The inverse \(T=4\ell_0-1\), \(u=\ell_1+jT\) proves bijectivity once the
prime, shell label, channel, orientation, and divisor support are retained.
The \(\varepsilon=0\) channel lies over terminal determinant parameter \(1\),
and the \(\varepsilon=1\) channel lies over \(p\).  The proof also derives the
complete selector from the factorization
\((R_jy-pa_j)(R_jz-pa_j)=p^2a_j^2\).  This is a structural equality of
integer presentations, not a map of the underlying geometric spaces; see
\cref{sec:es-s6-crosswalk}.

\item \Verified\ \textbf{The residue lattice and shell-count maps retain
their different domains.}
The lattice has the exact index-$72$ glue and the full coordinate
splitting in Theorems~\ref{esn:thm:niemeier}
and~\ref{esn:fa:thm:full-quotient}. Its $D_4$ shortest dual-class
sets are the vector, even half-spin and odd half-spin weights of the
flag tangent modules, in that marking
(Theorem~\ref{esn:thm:d4-f4-triality}).
For $p\equiv1\pmod{24}$ prime and $R>0$, $R\equiv3\pmod4$, put
$a=(p+R)/4$, $g=\gcd(R,pa)$, $R_0=R/g$, $N_0=pa/g$. Then
\[
 \left\{d>0:d\mid N_0^2,\ d\equiv-N_0\pmod{R_0}\right\}
 \xrightarrow{\ \cong\ }
 \left\{(b,c)\in\mathbb Z_{>0}^2:
                    \frac4p=\frac1a+\frac1b+\frac1c\right\},
\]
\[
 d\longmapsto
 \left(\frac{N_0+d}{R_0},\frac{N_0+N_0^2/d}{R_0}\right),
 \qquad U_R(p)=\frac{A_R(p)+\mathbf1_{R_0\mid2}}2.
\]
The inverse is $d=R_0b-N_0$. When $\gcd(R,pa)=1$ and $p\nmid a$,
the exact raw count is
$A_R(p)=c_R(a;-1)+2c_R(a;-p^{-1})$ in the exponent-box notation
of Theorem~\ref{esn:thm:raw-channels}. Equal prime residues do not
determine these counts (Theorem~\ref{esn:thm:non-descent}).
The finite-sample map in Theorem~\ref{esn:thm:sample-fourier}
sends these counts to an integral augmentation vector. Its Fourier
inversion recovers the aggregate residue-count function; the centering
$Q_{R,P}=6W_{R,P}-M_{R,P}u$ omits its constant Fourier component.
The sample $P$ remains in the domain, and no universal selector
nonemptiness is inferred.

\item \Verified\ \textbf{Two complex columns give the exact
four-dimensional \(SU(3)\) Gram cone.}
\[
 \operatorname{Mat}_{3\times2}(\mathbb C)/SU(3)
 \cong\operatorname{Herm}_2(\mathbb C)_{\geq0}.
\]
For the one-point compactified column representation,
\[
 S^{12}_{\rm col}/SU(3)\cong B^4,
\]
and the rank-one, zero, and infinity orbit types form its boundary \(S^3\).
This representation sphere is not the cusp-smash \(S^{12}\); see
\cref{thm:hrgc-gram-quotient,thm:hrgc-s12-ball-quotient}.

\item \Verified\ \textbf{Three complex columns give an exact cubic
rank-three cone with a distinguished rank-loss \(S^8\).}
For \(\mathscr V_3=\operatorname{Mat}_{3\times3}(\mathbb C)\), the complete
orbit invariants are
\[
 \mathscr V_3/U(3)\cong\operatorname{Herm}_3(\mathbb C)_{\geq0},
 \qquad
 \mathscr V_3/SU(3)
 \cong\{(H,z):H\geq0,\ |z|^2=\det H\}.
\]
The second space is \(\mathbb R^{10}\), so
\[
 S^{17}/SU(3)\cong S^9,\qquad
 \mathscr V_3^+/SU(3)\cong S^{10}.
\]
Inside that \(S^{10}\), the compactified rank-deficient locus is the
standard topologically locally flat \(S^8\), with
\[
 S^{10}\setminus S^8\cong\mathbb R^8\times\mathbb C^*\simeq S^1.
\]
The generator is the determinant phase \(\det A/|\det A|\), and the
trace-one rank-one stratum is \(\mathbb CP^2\).  Quotienting the determinant
phase gives \(\mathscr V_3^+/U(3)\cong B^9\) with boundary \(S^8\).  As a
whole this is a cubic Jordan/Gram carrier rather than a rank-two quadratic
Lorentz cone.  The marked angular quotient nevertheless gives an exact split
Lorentz observation with every complementary direction retained, as recorded
in the next item; see \cref{sec:hr3c-main}.

\item \Verified\ \textbf{The marked angular epimorphism gives a split
positive Lorentz observation, and its full fibre is retained.}
For
\[
 P=\begin{pmatrix}1&-1&0\\0&1&-1\end{pmatrix},\qquad
 u=(1,1,1)^{\mathsf T},\qquad
 R=P^*(PP^*)^{-1}
 =\frac13\begin{pmatrix}2&1\\-1&1\\-1&-2\end{pmatrix},
\]
put \(\Pi=uu^*/3\), \(Q_u=I_3-\Pi\), and
\[
 \mathcal O_P(H)=PHP^*,\qquad \jmath_P(G)=RGR^*.
\]
These positive maps satisfy
\[
 \mathcal O_P\jmath_P=\operatorname{id},\qquad
 \jmath_P\mathcal O_P(H)=Q_uHQ_u,
\]
with
\[
 \ker\mathcal O_P
 =\{a\,uu^*+uv^*+vu^*:a\in\mathbb R,\ v\in u^\perp\},\qquad
 \mathcal O_P^{-1}(0)\cap\mathcal P_3
 =\{\lambda uu^*: \lambda\geq0\}.
\]
Writing \(X=(R\;u)\), every Hermitian matrix is uniquely
\[
 H=X\begin{pmatrix}G&\psi\\\psi^*&\lambda\end{pmatrix}X^*,\qquad
 (G,\psi,\lambda)=
 \left(PHP^*,\frac13PHu,\frac19u^*Hu\right),
\]
so
\[
 \operatorname{Herm}_3(\mathbb C)
 \cong\operatorname{Herm}_2(\mathbb C)\oplus\mathbb C^2\oplus\mathbb R,
 \qquad 9=4+4+1,
\]
\[
 \det H=\lambda\det G-\psi^*\operatorname{adj}(G)\psi,\qquad
 H\succeq0\Longleftrightarrow
 G\succeq0, \psi\in\operatorname{im}G,
 \ \lambda\geq\psi^*G^+\psi.
\]
The faithful embedding
\[
 T_P(g)=RgP+\Pi=X(g\oplus1)X^{-1}
\]
acts by
\((G,\psi,\lambda)\mapsto(gGg^*,g\psi,\lambda)\).  For
\(g\in SL_2(\mathbb C)\), the summands are a Lorentz four-vector, the
defining complex spin representation, and a real scalar.  The vector action
factors through \(SO^+(1,3)\), while \(-I_2\) negates \(\psi\), so the full
carrier retains the spin double cover.

The section is proper and compactifies to
\[
 \jmath_P^+:B^4\cong\mathcal P_2^+
 \lhook\joinrel\longrightarrow(\partial\mathcal P_3)^+\cong S^8,\qquad
 e_2(\jmath_P(G))=\frac13\det G,
\]
and
\[
 \det(uu^*+\varepsilon\jmath_P(G))=\varepsilon^2\det G.
\]
The reverse boundary observation is nonproper.  In the strict \(SU(3)\)
quotient,
\[
 |\omega|^2=\lambda\det G-\psi^*\operatorname{adj}(G)\psi,
\]
and \(\omega=0\) is the rank-loss cone compactifying to
\(S^8\subset S^{10}\).  The four-branch chart supplies the exact toric seed
\[
 K(z)=\operatorname{diag}(z_0z_3,z_1z_2),\qquad
 \det K=z_0z_1z_2z_3,
\]
whose positive \(SL_2(\mathbb C)\)-orbit fills the future causal cone.
For \(Z\in\operatorname{Mat}_{3\times2}(\mathbb C)\) and retained
\(\alpha,c>0\),
\[
 p_\alpha(Z)=\alpha Z^*Z,\qquad
 m_{\alpha,c}^2(Z)=\frac{\alpha^2}{c^2}
 \sum_{a<b}|Z_{a1}Z_{b2}-Z_{a2}Z_{b1}|^2.
\]
This is the complete \(SU(3)\)-orbit momentum and is Lorentz covariant.
Its positive-mass locus is exactly rank two, and
\(m_{\alpha,c}^2(rZ)=r^4m_{\alpha,c}^2(Z)\), proving that the unrestricted
kinematic carrier has no positive threshold.  The exact reduction separates
four statements which the weak \(Z\)-equation alone does not.  First, for
\(m_0>0\) the colour quotient and reparametrization characteristic give
\[
 T^*\mathcal H^+_{m_0,c}
\]
with its canonical symplectic form, exactly the future spinless massive
particle.  Second, the intrinsic nonzero rank-one quotient gives
\(T^*\mathcal N_+^\times\).  Third, the full massless multiplier action is not
the standard massless variational system: the determinant constraint has zero
differential on rank one, its equation is \(\dot x=\lambda p\) without
\(\lambda=e\), four transverse rank-opening directions are non-gauge
first-order null directions, and \(Z=0\) admits arbitrary \(x(\tau)\).
Fourth, on the hard \(\omega_0\ne0\) packet the momentum differential has
rank four above threshold and rank three at threshold, so its displayed
action has \(\dot x=0\) in the strict interior and
\(\dot x\in\mathbb Rp\) only at the endpoint.  The packet's multiplication
operator is not the Hamiltonian generating this action; its positive endpoint
is a constraint-selected algebraic/free spectral threshold rather than a
dynamically generated interacting gap.  See
\cref{hlob:thm:split-observation,hlob:thm:block-coordinates,%
hlob:thm:lorentz-action,hlob:thm:cone-section,hlob:cor:Y3-block-form,%
hlob:thm:colour-momentum,wpr:thm:massive,wpr:thm:massless,%
wpr:thm:zero,wpr:thm:packet-action}.

\item \Verified\ \textbf{The determinant normal slice is an exact Lorentzian
host for explicit field models, but it does not itself select a mass gap.}
The real-linear isomorphism
\[
 \jmath_P:\operatorname{Herm}_2(\mathbb C)\xrightarrow{\ \sim\ }
 \mathcal N_P:=\{Y\in\operatorname{Herm}_3(\mathbb C):Yu=0\},
 \qquad G\longmapsto RGR^*,
\]
has inverse \(Y\mapsto PYP^*\), and
\[
 \det(uu^*+\varepsilon\jmath_P(G))=\varepsilon^2\det G.
\]
Writing
\(G(t,x,y,z)=
\bigl(\begin{smallmatrix}t+y&x+iz\\x-iz&t-y\end{smallmatrix}\bigr)\),
polarization gives
\[
 \eta_P(\jmath_P(G),\jmath_P(K))=tt'-xx'-yy'-zz'.
\]
The positive-semidefinite matrices are the future cone, the nonzero rank-one
matrices its null boundary, and the lifted congruence action induces
\(PSL_2(\mathbb C)\cong SO^+(1,3)\).  The future timelike determinant shell,
future-null locus, and candidate spacelike determinant shell are respectively
\[
 SL_2(\mathbb C)/SU(2),\qquad
 SL_2(\mathbb C)\cdot\operatorname{diag}(1,0),\qquad
 SL_2(\mathbb C)/SU(1,1).
\]

The standard metrics already retained by the Moore--Penrose construction
select
\[
 W_P=R^*R=C^{-1},\qquad
 K_P=\{a:a^*W_Pa=W_P,\ \det a=1\}\cong SU(2),
\]
and the positive form
\(E_P(G,K)=\operatorname{tr}(GW_PKW_P)\).  Under \(K_P\),
\[
 \operatorname{Herm}_2(\mathbb C)
 \cong\boldsymbol1\oplus\boldsymbol3,
\]
with trivial central action.  If this pointwise symmetry is promoted to a
gauge symmetry, the action
\[
 S[\mathcal A,G]=\int_M\left(
 \frac{|F_{\mathcal A}|^2}{2g_{\rm YM}^2}
 +\frac{Z_G}{2}|D_{\mathcal A}G|_{E_P}^2
 +\frac{m_G^2}{2}|G|_{E_P}^2
 +\frac{\lambda_4}{4}(\det G)^2\right)d\operatorname{vol}_g
\]
is nonnegative and gauge invariant, and
\(\widehat{\mathcal A}=R\mathcal A P\) intertwines the covariant
derivatives.  The normal field alone cannot screen a fundamental centre
charge; the full \((G,\psi,\lambda)\) carrier contains the centre-nontrivial
spinor \(\psi\), but its kinetic and interaction theory is not inserted by
this construction.

The candidate inverse-dual matrix and the positive Gram momentum are also
separated by an exact invariant:
\[
 \det G_{\eta,T}=-\frac23D_T<0,
 \qquad
 \operatorname{im}p_\alpha=
 \operatorname{Herm}_2(\mathbb C)_{\geq0}.
\]
Neither positive rescaling nor complex congruence can identify their inertia
types; a further morphism is required.  After a separate length/time
dictionary is supplied, the same affine host supports the ordinary
Klein--Gordon theory.  Its Fock Hamiltonian has excitation gap exactly
\(m_{\mathrm{KG}}\) for \(m_{\mathrm{KG}}>0\) and gap zero for
\(m_{\mathrm{KG}}=0\).  On a periodic spatial gauge regulator, the massless
transverse Gaussian frequency is \(2\pi/L\to0\).  Explicit countermodels also
show that entanglement occurs in both gapped and gapless systems.  A positive
auxiliary gap transfers to the physical Hamiltonian only through a map on the
entire excitation form domain satisfying
\[
 \|Ru\|\geq c_R\|u\|,\qquad
 \mathfrak b[Ru]\leq C_R
 (\mathfrak a_{\rm phys}[u]-E_{\rm vac}\|u\|^2),
\]
in which case
\(\operatorname{gap}(A_{\rm phys})\geq\gamma c_R^2/C_R\).
An intertwiner on a single gapped subspace is not enough.  See
\cref{hnsp-lorentz-theorem,hnsp-spin-theorem,%
hnsp-intrinsic-compact-representation,hnsp-intrinsic-gauge-scalar-theorem,%
hnsp-centre-screening-corollary,hnsp-candidate-spacelike,%
hnsp-determinant-orbit-geometry,hnsp-free-field-gap,%
hnsp-gauge-gaussian-gap,hnsp-coercive-transfer}.

\item \Verified\ \textbf{The marked cubic cone contains exact Lorentz
sections, but these do not identify it with the toroidal degeneration.}
With the oriented cross-product matrix \([a]_{\times}\),
\[
 \det\bigl(S+i[a]_{\times}\bigr)=\det S-a^{\mathsf T}Sa.
\]
Consequently each marked positive monodromy form \(B\) determines the exact
linear section
\[
 \ell_B(\lambda,a)=\lambda B+i[a]_{\times},\qquad
 \ell_B^{-1}\!\left(\operatorname{Herm}_3(\mathbb C)_{\geq0}\right)
 =\{\lambda\geq0:a^{\mathsf T}Ba\leq\lambda^2\det B\}.
\]
This is a future \((3+1)\)-Lorentz cone.  An explicit real-linear order-cone
map \(\Theta\) from \(2\times2\) Hermitian matrices satisfies
\[
 \det_3\Theta(H)=\frac{\operatorname{tr}H}{2}\det_2H
\]
and sends positive-semidefinite ranks \(2,1,0\) to \(3,2,0\).  It is not a
Jordan homomorphism: for a rank-one projection \(P\), the matrix
\(\Theta(P)\) has eigenvalues \(0,1/2,1\) and is not idempotent.  Neither map carries a toric period lattice, monodromy, or
singular-fibre filling.  In particular, the toric limit at fixed \(B=I_3\)
 is not the rank-drop operation \(\det B\to0\).  The rejected packet's claimed
index-three translation compression is likewise not a constructed map: the
actual six-to-two bridge is the saturated rank-lowering quotient
\(P:\mathbb Z^3\twoheadrightarrow\mathbb Z^2\), which sends the even and odd
Kuhn labels to the two unimodular target-triangle orbits.  See
\cref{thm:teb-hermitian-determinant-splitting,%
prop:teb-theta-order-lift,prop:teb-fan-morphism}.

\item \Verified\ \textbf{The rank-three toroidal successor maps onto the
candidate cusp with a connected elliptic kernel.}
The Kuhn fan gives a proper analytic family
\[
 \pi_3:N_3\longrightarrow\Delta,\qquad
 N_{3,t}\cong E_\tau^3\quad(t\ne0),
\]
whose deepest local equation is \(t=z_0z_1z_2z_3\), whose cusp matrix is
\(M_3=\bigl(\begin{smallmatrix}I_3&I_3\\0&I_3\end{smallmatrix}\bigr)\),
and whose filling map kills the three angular cycles and retains the three
deck cycles.  The irreducible central fibre has
\[
 (f_0,f_1,f_2,f_3)=(1,7,12,6),\qquad \chi(W_3)=6.
\]
Its normalization is the permutohedral threefold with orbit vector
\((1,14,36,24)\); the positive-codimension strata map with degrees
\(2,3,4\) onto the \(7,12,6\) quotient strata.  Hence any real
\(S^8\)-total-space fibration retaining \(W_3\) requires the remaining
singular fibres to contribute Euler sum \(-4\).
For an injective integral shear
\(\mathcal B:\bar L_3\hookrightarrow L_3\), this becomes the exact
determinant/index formula
\[
 \chi(W_{3,\mathcal B})
 =3!\,[L_3:\mathcal B\bar L_3]
 =3!\,|\det\mathcal B|.
\]

The explicit matrices
\[
 P=\begin{pmatrix}1&-1&0\\0&1&-1\end{pmatrix},
 \qquad
 Q=\begin{pmatrix}0&-1&1\\1&-1&0\end{pmatrix}
\]
define a surjection
\(\Phi=P\oplus Q:H_1(T^6)\twoheadrightarrow H_1(T^4)\) satisfying
\[
 \Phi M_3=M_F\Phi,\qquad c_F\Phi=Qc_3,\qquad
 PZ_3(t)=Z_2(t)Q.
\]
It extends over all cusp toric strata to a holomorphic map with connected
elliptic kernel.  Integrally,
\[
 M_3\cong M_F\oplus
 \begin{pmatrix}1&1\\0&1\end{pmatrix},
\]
and the order-three and order-four free affine peripheral actions have
compatible rank-six lifts.  The six Kuhn tetrahedron orbits map
\(S_3\)-equivariantly to the six marked diagonal Weyl chambers, but only as
fan/Cartan data; the four-branch toroidal germ and the quadratic Gram-rank
germ remain inequivalent.  What remains open is the single holomorphic
period-and-overlap family across both finite-order discs.  See
\cref{hrftd:cor:euler-shear-determinant,%
thm:teb-holomorphic-cusp-bridge,prop:teb-explicit-conjugacy,%
 prop:teb-finite-monodromy-lift,prop:teb-weyl-chambers}.

\item \Verified\ \textbf{The free rank-six cusp lattice, the six chambers,
and the Euler factor six are now separated exactly.}
The lattice native to the higher-rung construction is
\[
 \Lambda_{\rm tor}=H_1(T^6;\mathbb Z)
 =\mathbb Z^3_{\rm angular}\oplus\mathbb Z^3_{\rm deck}.
\]
For \(k=(1,1,1)\), the map \(\Phi=P\oplus Q\) gives an exact
\(S_3\)- and monodromy-equivariant sequence
\[
 0\longrightarrow\mathbb Z(k,0)\oplus\mathbb Z(0,k)
 \longrightarrow\Lambda_{\rm tor}
 \longrightarrow A_F\oplus B_F\longrightarrow0.
\]
It splits integrally as a monodromy module but not over
\(\mathbb Z[S_3]\).  The quotient copies are \(A_2\) weight lattices; their
root sublattices have index three.  The six standard homology generators are
two permutation orbits of size three, whereas the six Kuhn cells and six
marked Weyl chambers form one regular \(S_3\)-orbit.  Consequently
\[
 \operatorname{rk}\Lambda_{\rm tor}=6,\qquad
 \#\{\text{chambers}\}=3!=6,\qquad
 \chi(W_{3,\mathcal B})=3!|\det\mathcal B|
\]
do not identify six lattice directions with six strata.  The regular
six-element chamber set instead occurs in the rank-nine bilinear carrier
\[
 \operatorname{Hom}(\overline L_3,L_3)
 =\bigoplus_i\mathbb ZE_{ii}\oplus
  \bigoplus_{i\ne j}\mathbb ZE_{ij}:
\]
the off-diagonal units are a regular \(S_3\)-set, and
\(c_{(\pi_1\pi_2\pi_3)}\mapsto E_{\pi_1,\pi_3}\) is an equivariant
bijection.  These are exactly the mixed angular--deck basis pairs with zero
alternating pairing; the three excluded diagonal pairs have pairing one.
The ambient real matrix representation is equivariantly isomorphic to
\(\operatorname{Herm}_3(\mathbb C)\), but the symmetric/skew integral
splitting has index \(2^3\), and the images of the off-diagonal units are
indefinite.  This is a representation/incidence bridge, not a positive,
Jordan, determinant, or toric morphism.  The actual
nine-dimensional-boundary map is only the piecewise-linear angular map
\(\Phi_{\rm G}\circ\jmath\circ P\); the deck lattice is instead used by the
period and filling squares.  This map does not descend to a torus fibre.
The separate bundle-defect lattice \(H^4(U;\mathbb Z)\cong(\mathbb Z^2)^3\)
has a different domain and degree.  See
\cref{rsl:sec:rank-six,rsl:thm:three-sixes,rsl:prop:gram-scope,%
rnmt:thm:regular-torsor,rnmt:prop:representation-bridge}.

The successful lifted candidate actions are distinct from the diagonal
local-system candidate
\[
 M_1=I_3\otimes S,\qquad M_2=I_3\otimes Q,\qquad M_0=I_3\otimes T,
 \qquad M_1M_2M_0=I.
\]
Its finite factors have orders \(4\) and \(6\), but
\(I-M_j:T^6\to T^6\) is an isogeny for \(j=1,2\); every affine lift of
either finite action therefore has a fixed point.

\item \Refuted\ \textbf{There is no literal smooth complex-fourfold
analogue on a homotopy \(S^8\).}
An almost-complex tangent bundle would have
\[
 c_1=c_2=c_3=0,\qquad c_4=2u,\qquad p_2=2c_4=4u,
\]
whereas the signature formula forces \(p_2=0\).  A nonsingular
\(T^6\)-bundle over \(S^2\) with total space \(S^8\) is independently
excluded by both the homotopy sequence and Euler characteristic.  A real
singular \(T^6\)-over-\(S^2\) construction remains a legitimate search
problem, but its special fibres must obey
\(\sum_i\chi(F_i)=2\).  For monodromy \(\rho\) on the regular piece, its
actual filling inclusions must use the corrected Serre filtration
\[
 0\longrightarrow(\textstyle\bigwedge^q\Lambda)_\rho
 \longrightarrow H_q(E;\mathbb Z)
 \longrightarrow
 H_1(B^\circ;\textstyle\bigwedge^{q-1}\Lambda)
 \longrightarrow0,
\]
and the dual quotient
\(H^q(E;\mathbb Z)\twoheadrightarrow
(\bigwedge^q\Lambda^\vee)^\rho\).  Thus the integral Mayer--Vietoris test
uses coinvariants as the homological subobject and the base-one-cycle term
as quotient, not conversely; see
\cref{hrs8:sec:constraints,hrs8:eq:serre-regular-part,%
 hrs8:eq:serre-cohomology-invariants}.  Every effective torus action on
\(S^8\) has rank at most four, and the sharp coordinate model is
\[
 S^8/T^4\cong\operatorname{Susp}(\Delta^3)\cong B^4.
\]
This is a real two-pole singular compactification with generic \(T^4\)
orbit, not a three-special-fibre \(T^6\)-fibration over \(S^2\); see
\cref{hrs8:eq:torus-rank,hrs8:eq:t4-quotient}.

\item \Verified\ \textbf{The regular locus of the reconstructed \(S^6\)
has a six-coordinate, nonextendable defect lattice.}
For \(D=W\sqcup S_1\sqcup S_2\) and \(U=X\setminus D\),
\[
 \boxed{H^4(U;\mathbb Z)
 \cong H_1(W)\oplus H_1(S_1)\oplus H_1(S_2)
 \cong\mathbb Z^2\oplus\mathbb Z^2\oplus\mathbb Z^2.}
\]
Principal \(SU(3)\)-bundles on \(U\) are classified exactly by this
\(c_2\)-class, and the pair connecting morphism
\(H^4(U)\xrightarrow{\sim}H^5(X,U)\) proves that extension across all
three special fibres occurs exactly for \(c_2=0\).  The displayed group is
the free rank-six lattice \(\mathbb Z^6\), not a cyclic \(\mathbb Z/6\).

\item \Verified\ \textbf{The primitive fibre degree has exact order-twelve
and order-six normal-holonomy quotients.}  Complete proofs, including the
one-shot collar matrix, transpose system, saturated kernel, coefficient
reduction, tensor-centre map, and common-centre quotient, are given in
\cref{zdc:sec:bridge,zdc:thm:fibre-degree,zdc:prop:finite-reductions,%
zdc:thm:tensor-centres,zdc:thm:common-centre}.  Those group and lattice
theorems do not themselves define a gauge field or an energy operator.

In the marked basis
\[
 (a_0,b_0,a_1,b_1,a_2,b_2)
 =(\bar\gamma,\bar u,g_1,h_1,g_2,h_2),
\]
the transpose of the actual one-shot collar matrix gives
\[
 \boxed{[\deg_F]=(12,0,4,0,-3,0)},
\]
meaning
\[
 \operatorname{res}_F\!\left(\sum_{i=0}^2(n_i a_i+r_i b_i)\right)
 =(12n_0+4n_1-3n_2)\omega_F.
\]
It is primitive, with saturated kernel
\[
 \boxed{\ker\deg_F
 =\mathbb Z\langle
 b_0,b_1,b_2,a_0-3a_1,3a_1+4a_2\rangle,}
\]
Define
\[
 q_{\rm col}:L\longrightarrow H_1(X^\circ;\mathbb Z),\qquad
 a_0\mapsto c,\quad a_1\mapsto\bar g_1,\quad a_2\mapsto\bar g_2,\quad
 b_0,b_1,b_2\mapsto0.
\]
If \(\Phi_1:H_1(X^\circ;\mathbb Z)\to\mathbb Z\) has row
\((12,4\ell_1,3\ell_2)\) in the basis \((c,\bar g_1,\bar g_2)\), then the
factorization through the same collar functional that kills the final cusp is
\[
 \deg_F=\Phi_1\circ q_{\rm col},\qquad
 \Phi_1(j_*[\widetilde a_0])=p.
\]

Let \(\vartheta_1(a_1)=e^{-2\pi i/3}\) and
\(\vartheta_2(a_2)=e^{-2\pi i/4}\) be the local normal-line holonomies,
trivial on the other marked generators.  The actual relative twist gives
\[
 \chi_{12}\!\left(\sum_{i=0}^2(n_i a_i+r_i b_i)\right)
 =\vartheta_1(n_1a_1+r_1b_1)
  \vartheta_2(n_2a_2+r_2b_2)^{-1}
 =e^{-2\pi i(4n_1-3n_2)/12}.
\]
If \(\rho_m(k)=e^{-2\pi ik/m}\), then
\[
 \boxed{\chi_{12}=\rho_{12}\circ\deg_F,\qquad
        \chi_6:=\chi_{12}^2=\rho_6\circ\deg_F.}
\]
Therefore
\[
 H^4(U;\mathbb Z)/\ker\chi_{12}\cong\mathbb Z/12\mathbb Z,
 \qquad
 H^4(U;\mathbb Z)/\ker\chi_6\cong\mathbb Z/6\mathbb Z.
\]
Ordinary coefficient reduction instead gives
\[
 H^4(U;\mathbb Z/6)\cong(\mathbb Z/6)^6;
\]
it is not the single distinguished cyclic quotient selected by \(\chi_6\).
For a principal \(SU(3)\)-bundle \(P\to U\), the order-six character is
exactly
\[
 \chi_6(c_2(P))
 =\exp\!\left(-\frac{2\pi i}{6}
 \langle c_2(P)|_F,[F]\rangle\right).
\]
This is a topological congruence for the fibre second-Chern number, not a
line-operator or spectral assertion.

The tensor embedding
\[
 SU(2)\times SU(3)\hookrightarrow SU(6),\qquad
 (A,B)\longmapsto A\otimes B
\]
restricts to an isomorphism
\(\mu_2\times\mu_3\xrightarrow{\sim}Z(SU(6))=\mu_6\).  Since \(U\) has
homotopy dimension at most five, \(c_2\) classifies \(SU(n)\)-bundles for
\(n\ge3\), while
\[
 c_2(E_2\otimes E_3)=3c_2(E_2)+2c_2(E_3).
\]
Every chosen \(SU(6)\)-bundle class \(h\in H^4(U;\mathbb Z)\) therefore
admits a tensor-subgroup reduction, and all its \(c_2\)-factorizations are
\[
 (e,f)=(h+2t,-h-3t),\qquad t\in H^4(U;\mathbb Z).
\]
No nonzero \(h\) or reduction is canonically selected.  Separately,
\[
 H^5(U;\mathbb Z)
 \cong\ker\{\mathbb Z^3\xrightarrow{(1,1,1)}\mathbb Z\}
 \cong Q(A_2),
\]
for the combinatorial form making the three distinguished support basis
vectors orthonormal.  This is not a cup-product or intersection form, and
the support-redistribution lattice does not identify the rank-five kernel of
\(\deg_F\) with \(A_5\).

The finite-group isomorphism
\(H^4(U)/\ker\chi_6\cong Z(SU(6))\) supplies no \(SU(6)\) transition
functions or connection.  The injective tensor map does not factor through
\((SU(2)\times SU(3))/\mu_6\), because the central \(\mu_6\) maps to
nontrivial scalar matrices.  The distinct common-centre map is
\[
\begin{aligned}
 \Gamma_6
 &=\bigl\{((\lambda\mu)^{-1},\lambda I_2,\mu I_3):
             \lambda\in\mu_2,\ \mu\in\mu_3\bigr\},\\
 \frac{U(1)\times SU(2)\times SU(3)}{\Gamma_6}
 &\lhook\joinrel\longrightarrow U(6),\\
 [z,A,B]&\longmapsto z(A\otimes B),
\end{aligned}
\]
where \(\Gamma_6\cong\mu_6\).  Its determinant is \(z^6\), and its image
meets \(SU(6)\) in exactly \(\iota(SU(2)\times SU(3))\); the image is not
generally contained in \(SU(6)\).  The missing step is a degree-changing
identification with a degree-two common-centre lifting obstruction and its
line-operator pairing.  The separate global clutching classes on \(S^6\)
have \(c_3=\pm2ku\), \(k\in\pi_5(SU(3))\cong\mathbb Z\), and restrict
trivially to \(U\); see \cref{hrs8:subsec:su3-regular-locus}.

\item \Verified\ \textbf{The exact typed transfer from auxiliary
Rayleigh quotients to a physical excitation spectrum is proved.}
Let \(H\ge0\) be self-adjoint on \(\mathcal H_{\mathrm{phys}}\), and set
\[
 P_0=\mathbf1_{\{0\}}(H),\qquad
 \mathcal H_{\mathrm{exc}}=(I-P_0)\mathcal H_{\mathrm{phys}},\qquad
 \mathfrak h[\psi]=\lVert H^{1/2}\psi\rVert^2,\qquad
 \mathcal Q(H)=D(H^{1/2}).
\]
For a nonnegative quadratic form \(\mathfrak q_n\) and nonzero
\(u_n\in D(\mathfrak q_n)\subset\mathcal K_n\), suppose
\[
 J_n:\operatorname{span}\{u_n\}\longrightarrow\mathcal Q(H),\qquad
 \psi_n=(I-P_0)J_nu_n,
\]
and, for \(c>0\), \(C_n,\varepsilon_n\ge0\),
\[
 \lVert\psi_n\rVert\ge c\lVert u_n\rVert,\qquad
 \mathfrak h[\psi_n]\le C_n\mathfrak q_n[u_n]
 +\varepsilon_n\lVert u_n\rVert^2,\qquad
 C_n\frac{\mathfrak q_n[u_n]}{\lVert u_n\rVert^2}
 +\varepsilon_n\longrightarrow0.
\]
Then the self-adjoint excitation restriction
\[
 H_{\mathrm{exc}}
 =H\big|_{D(H)\cap\mathcal H_{\mathrm{exc}}}
\]
satisfies
\[
 \inf\sigma(H_{\mathrm{exc}})=0,\qquad
 0\in\sigma_{\mathrm{ess}}(H_{\mathrm{exc}}).
\]
The projection removes the full kernel of \(H\); orthogonality to a
single chosen vacuum is equivalent only after vacuum uniqueness is separately
proved.  This is a proved implication with fully stated antecedents, not a
claim that the present geometry supplies them.  See
\cref{cusp-interpretation-typed-physical-transfer}.

\item \Open\ \textbf{The physical Yang--Mills bridge data are absent.}
The workbench has not constructed a common gauge-invariant positive-metric
\(\mathcal H_{\mathrm{phys}}\), a four-dimensional quantum Yang--Mills
Hamiltonian \(H\), or maps \(J_n\) obeying the norm and energy inequalities
through gauge reduction, reflection-positive reconstruction,
renormalization, and the continuum and infinite-volume limits.  Therefore
the verified radial scalar spectra, auxiliary direct integral, classical
gauge Hessians, finite cochain operators, compact-link flux configurations,
 finite Gaussian oscillators, magnetic Laplacians, and reducible
 split-Cartan \(SU(r)\) Jacobi operators do not prove or disprove a mass gap for a physical quantum
Yang--Mills Hamiltonian.  See
\cref{cusp-gauge-missing-physical-bridge,%
cusp-interpretation-present-boundary}.

\end{enumerate}

\paragraph{Exact-computation certificates.}
The retained exact-arithmetic harnesses passed all \(28\) full-manuscript
checks, all \(38\) independently implemented core checks, and all \(13\)
cusp-spectral checks, with no failures; the \(13\)-check harness also passed
unchanged under \texttt{python -O}.  The first two harnesses cover finite and
cusp monodromy, invariant and coinvariant
lattices in exterior degrees \(0\) through \(4\), Smith forms and
saturation, finite-quotient norm operators, the relation-group determinant,
specialization indices, three twisting parameters, period-form identities,
and the corrected existential form of Remark~7.28.  The \(13\)-check
cusp-spectral harness is scoped to the radial period determinant and full
inverse transpose, the restricted dual maps and pairings, the integral
coframe wedge identity, the ambient Koide discriminant identity, the finite
coordinate eigenvalues, and the radial volume and first-correction
expansions.  The geometric collapse, finite-band and cusp inclusion maps,
Leray comparison, Kummer morphisms, cone carriers, compact-link construction,
and Koide counting theorem are proved separately rather than inferred from
those executions.

\paragraph{Independent review checkpoint.}
Reviewer~B independently reviewed the reconstruction existing up to its
evidentiary cutoff.  It withdrew the earlier topology-based rejection,
distinguished the intrinsic specialization \(\operatorname{sp}^q\) from the
then-open collapse pullback \(c^*\), and found no load-bearing defect in the
non-collapse nearby-cycle/Leray and based van Kampen route.  Its
provisionally favourable disposition and focused verification requests are
useful independent evidence about the argument it saw.  The later Codex
proofs in this workbench are separate and should not be attributed to that
review.  In its 26 August analytic follow-up, Reviewer~B independently proved
the finite-complex/UCT top-degree base-change statement on every
scheme-theoretic surface fibre and, from CAND-X, CAND-W, and TOP-H2, confirmed
the conductor--duality implication
$R^2f_*(T_X\otimes L)\ne0$ for every $L$ \cite{ReviewerBValidation2026}.
Those three inputs remain hypotheses in that review record; its cutoff labels
do not replace their later workbench classifications.

\paragraph{Current evidence boundary.}
No mathematical gap is presently recorded in the construction, either
topology certificate, sphere recognition, algebraic dimension, or Hodge
and Bott--Chern/Aeppli calculations.  Exact marked-fibre period-coordinate
representatives, global Bott--Chern generators, and their complete product
structure are now included.  The canonical data do not single out a literal
global coefficient atlas for the chosen smooth representatives through the
singular fillings; no such gauge is needed for the class or algebra theorem.
The new cusp spectral statements apply to the explicitly declared quotient
metrics and auxiliary or classical operators.  They now include the complete
fixed-energy lattice-point law
\(N_T(E)=\pi D_T\sum_n(E/(4\pi^2)-\lVert n\rVert^2)_++O_E(T+1)\)
and the full-torus integral magnetic bound
\(\operatorname{gap}(\Delta_{\mathsf E,T})
\le4\pi\sqrt{|\operatorname{Pf}(\mathsf E)|/D_T}\).
The modified-action toric Kummer map now completes the local
degree-\(N^2\) regular-fibre cover, while
retaining its direct degree \(N^3\), ramification, nonflat relative map, and
strictly local scope.  The exact maps \(q_2^{\wedge2}\),
\(q_{1,\mathrm{cusp}}^{\wedge4}\), and \(q_1^{\wedge4}\) now provide
topological \(S^{12}\)- and \(S^{24}\)-source carriers only by changing the
targets through suspension, collapse, and coning.  Two copies are rank-minimal
cusp-locally and four are rank-minimal over the full monodromy base.  Further
quotients \(\widehat q_3^{\wedge2}:S^{12}\twoheadrightarrow S^6\) and
\(\widehat q_\gamma^{\wedge4}:S^{24}\twoheadrightarrow S^8\) contain moving,
null-homotopic metric four-spheres with eventual exact Euclidean basepoint
balls.  Both full maps are themselves based null-homotopic.  More generally,
\(\widehat q_\gamma^{\wedge r}\) and
\(\widehat q_3^{\wedge r}\) are based null-homotopic for every \(r\ge2\).
The one-copy class of \(\widehat q_3\) is zero by the explicit framed
filling of Theorem~\ref{h24:thm:cusp-g-zero}; that of
\(\widehat q_\gamma\) remains open.
The conical-target classes of \(q_2^{\wedge2}\),
\(q_{1,\mathrm{cusp}}^{\wedge4}\), and \(q_1^{\wedge4}\) remain \Open.  The
topological five-manifolds \(\mathfrak E_{12},\mathfrak E_{24}\) realize
\(\mathbb R^4\) as a decompactification boundary fibre.  Unrescaled slices
collapse in every fixed compatible ambient metric; for the fixed
cusp-adapted metrics, the rescaled slice bubbles are \(\mathbb R^4\) while
the full ambient bubbles are \(\mathbb R^8\) and \(\mathbb R^{12}\).
This leaves intact the
nonextension to the fixed product target and supplies no compact complex
structure on either smash sphere.  The repaired compact-link calculation
keeps cover degree and mesh refinement independent and proves the exact
finite Gaussian gap and its zero-gap diagonal sequence, but not a
nonperturbative compact-link spectrum.  The physical Hilbert space, quantum
Hamiltonian, and comparison morphisms required by the typed transfer theorem
remain \Open; accordingly, these calculations neither prove nor disprove a
four-dimensional quantum Yang--Mills mass gap.  Separately, the exact
three-column quotient is now a rank-three cubic carrier with compactified
rank-loss pair \((S^{10},S^8)\), while a literal almost-complex \(S^8\) is
refuted.  The toroidal angular map gives that cubic carrier a split positive
observation onto the ordinary future Lorentz cone, with exact
\(9=4+4+1\) four-vector--spin-representation--scalar coordinates and a
proper \(B^4\hookrightarrow S^8\) section.  Its determinant-phase equation,
toric seed, and \(SU(3)\)-invariant momentum are explicit, and
\(m_{\alpha,c}^2(rZ)=r^4m_{\alpha,c}^2(Z)\) proves that the unrestricted
kinematic carrier has no positive threshold.  On the candidate's own regular
locus, the complete \(SU(3)\)-bundle
 classification is the free defect lattice
 \(H^4(U;\mathbb Z)\cong\mathbb Z^2\oplus\mathbb Z^2\oplus\mathbb Z^2\),
 and every nonzero class fails to extend across the three special fibres.
 The one-shot collar matrix computes the primitive fibre-degree row
 \((12,0,4,0,-3,0)\); the candidate's local normal holonomies are exactly its
 reductions \(\chi_{12}=\rho_{12}\circ\deg_F\) and
 \(\chi_6=\rho_6\circ\deg_F\).  The latter quotient is abstractly the same
 \(\mu_6\) as the centre of the tensor subgroup
 \(SU(2)\times SU(3)\hookrightarrow SU(6)\), but no degree-changing map to a
 common-centre lifting obstruction, gauge cocycle, connection, or physical
 line pairing has been constructed.
 The carrier theorem advances the physical bridge through an exact
colour-orbit momentum and a now fully reduced worldline action: its massive
stratum is the standard spinless particle, its punctured massless
characteristic quotient is standard while its multiplier action and apex are
irregular, and the hard-constraint packet propagates only on its threshold
stratum.  Its normal-slice
continuation now supplies a constant affine Minkowski metric, the spin and
Poincar\'e actions, an exact Klein--Gordon/Fock theory and finite regulator,
an independently chosen classical compact-group Yang--Mills theory, and the
metric-selected \(K_P\cong SU(2)\) gauge--scalar model.  These constructions
show both possible gap values: the free scalar gap is the externally supplied
\(m_{\mathrm{KG}}\), while the massless scalar and thermodynamic Gaussian
gauge limits are gapless.  They also isolate the candidate cusp matrix as
spacelike and outside the positive Gram-momentum image.  Dynamics on the
retained \(\psi,\lambda,\omega\) directions, a dynamically selected
gauge-invariant interacting spectral sector, a continuum quantum
reconstruction, and a full-domain comparison map to its physical excitation
space remain unconstructed.  The labels above mean that
complete arguments are present here for audit.  Specialist
replication and external adjudication are
separate validation boundaries, neither hypotheses nor gates on these
proofs; they remain welcome.  The corrected CDP sufficient theorem is now
proved in
\cref{thm:corrected-cdp-global,thm:corrected-cdp-fibrewise}.  Recovery of
the original theorem under its original hypotheses would instead require
overturning a verified premise of the constructed counterexample.
